\PassOptionsToPackage{obeyspaces,spaces}{url}
\documentclass[11pt]{article}
\usepackage[T1]{fontenc}
\usepackage{lmodern}
\usepackage[margin=1in]{geometry}
\usepackage{amsmath,amssymb}
\usepackage{booktabs,longtable,array,tabularx}
\usepackage{enumitem}
\usepackage[dvipsnames]{xcolor}
\usepackage{url}
\usepackage[colorlinks=true,linkcolor=NavyBlue,citecolor=NavyBlue,urlcolor=NavyBlue]{hyperref}
\usepackage{microtype}
\hypersetup{pdftitle={Fortuitous states of N=4 SYM at weak coupling: results, code and cross-checks},
  pdfauthor={The Fortuity collaboration}}

\setlist{itemsep=2pt,topsep=3pt}
\setlength{\emergencystretch}{3em}
\renewcommand{\arraystretch}{1.12}
\urlstyle{tt}
\newcommand{\file}[1]{\path{#1}}
\newcommand{\R}[1]{\textsf{R-0#1}}
\newcommand{\T}[1]{\textsf{T-0#1}}
\newcommand{\SU}{\mathrm{SU}}
\newcommand{\SO}{\mathrm{SO}}
\newcommand{\USp}{\mathrm{USp}}
\newcommand{\Gtwo}{G_2}
\newcommand{\cN}{\mathcal{N}}
\newcommand{\ph}{\phi}
\newcommand{\ps}{\psi}
\newcommand{\ff}{f}
\newcommand{\arxiv}[1]{\href{https://arxiv.org/abs/#1}{arXiv:#1}}
\newcommand{\var}[1]{\texttt{\$\{#1\}}}

\title{Fortuitous states of $\cN=4$ SYM at weak coupling:\\[2pt]
\large results, code and cross-checks of the Fortuity collaboration}
\author{The Fortuity collaboration}
\date{14 October 2026}

\begin{document}
\maketitle

\noindent\fbox{\parbox{\dimexpr\textwidth-2\fboxsep-2\fboxrule}{\small
\textbf{How this document and the results in it were produced.} The computations, the code and this text were
produced by Claude, Anthropic's AI model, working as AI agents for the members of the Fortuity collaboration. The
members set the questions, read and checked the outputs, and reviewed one another's results; a separate set of AI
agents reviewed the results independently (section~\ref{sec:checks}). Every number below either carries an automated
certificate of the kind described in section~\ref{sec:method} or is marked as unverified; none has been rederived by
hand. The labels \textsf{R-0xx} (results) and \textsf{T-0xx} (computations) are those of the collaboration's records
and are used in the comments of the code; appendix~\ref{app:index} lists each of them with the code, the data and the
checks in this repository. This document supersedes the summary of 7 October 2026. Corrections are welcome.}}

\begin{abstract}
We summarise what the collaboration has established about the $\frac1{16}$-BPS cohomology of $\cN=4$ super
Yang--Mills theory at finite rank: the classical $Q$-cohomology in the BMN sector for $\SU(2)$ through level $j=44$,
$\SU(3)$ through $j=32$ and in one column at $j=34$, $\SU(4)$ through $j=30$, $\SO(5)$, $\SO(7)$, $\SO(9)$,
$\SO(11)$, $\USp(6)$, $\USp(8)$, $\USp(10)$ and $\Gtwo$; BMN indices for $\SU(5)$ and for $\SO(2N+1)$ against
$\USp(2N)$ up to $N=6$; the full-letter cohomology of $\SO(7)$ and $\USp(6)$ at levels 18 to 22; and the
perturbative fate of the lowest classes at one and two loops. The main results are:
(i) the $\SU(3)$ threshold operator of Choi, Choi, Kim, Lee and Lee is $Q$-exact, and the eigenvalue criterion on
which their non-graviton index rests is false; the first $\SU(3)$ non-graviton class is at $j=30$
($E=\tfrac{25}{2}$, $[3,0]$), it is the brane-fused operator of de Mello Koch, Kim, Kim, Lee and Lee, it is odd under
charge conjugation, and it survives one loop; every $\SU(3)$ black-hole class below $j=34$ is fused on an $\SU(2)$
black hole by scalar branes, and at $j=34$ the $[3,1]$ column contains, besides a descendant of the fused class at
$j=32$, exactly one class that is not of that form;
(ii) $\SU(2)$ has 20 non-graviton classes through $j=44$, of which 5 lift at one loop; the threshold class has no
admissible lifting partner at any order given a superconformal selection rule, $d_2$ vanishes on it by direct
computation, and the index alone protects it at every coupling;
(iii) $\SU(4)$ has no non-graviton class through $j=30$;
(iv) $\SO(2N+1)$ and $\USp(2N)$ first differ at $j=4N+6$ in the representation $\mathrm{Sym}^{N-3}[2,0]$ ($N=3,\dots,6$);
for $N=3,4,5$ the difference is one fermionic fortuitous $\SO(2N+1)$ multiplet per irrep, the lowest-level
fortuitous class of that group, the $\SO(7)$ class dressed by $N-3$ separated branes, and lifted at one loop;
(v) structural results: a charge-conjugation parity that makes every multi-graviton even, twisted and kernel-parity
indices that detect black holes without a cohomology computation, multi-graviton classes invisible on the Coulomb
branch, and a two-loop supercharge on derivative-free words fixed by nilpotency.
We compare every overlapping statement with the published literature, list the cross-checks, and say what remains
open. All code and data behind the results are in this repository.
\end{abstract}

\tableofcontents

\section{Main results at a glance}\label{sec:glance}

Table~\ref{tab:thresholds} collects the lowest non-graviton classes found in the BMN sector. ``Level'' is the power
$j$ of $t$ in the BMN index (section~\ref{sec:setup}); a complete scan of all levels $j<j_0$ certifies the lowest
\emph{level}, while the lowest \emph{energy} needs every level $j<3E$, since $E\ge j/3$ in this sector.

\begin{table}[h]\centering\footnotesize
\setlength{\tabcolsep}{4pt}
\begin{tabularx}{\textwidth}{@{}lllllX@{}}
\toprule
group & lowest non-graviton class & $E$ & complete scans & certified & quantum fate\\
\midrule
$\SU(2)$ & $(5;4,4,4)$, $[0,0]$, $j=24$ & $\tfrac{19}{2}$ & $j\le44$ & level, energy & survives 1 loop; all orders$^{a}$; index$^{b}$\\
$\SU(3)$ & $(5;7,4,4)$, $[3,0]$, $j=30$ & $\tfrac{25}{2}$ & $j\le32$ & level & survives 1 loop; orders 3, 4$^{c}$\\
$\SU(4)$ & none for $j\le30$ & -- & $j\le30$ & level $\ge32$ & --\\
$\SO(5)$ & none for $j\le28$ & -- & $j\le28$ & level $\ge30$ & --\\
$\SO(7)$ & $(1;3,3,3)$, $[0,0]$, $j=18$ & $\tfrac{17}{2}$ & $j\le28$ & level, energy & lifted at 1 loop\\
$\SO(9)$ & $(1;5,3,3)$, $[2,0]$, $j=22$ & $\tfrac{21}{2}$ & $j\le22$ & level & lifted at 1 loop\\
$\SO(11)$ & $(1;7,3,3)$, $[4,0]$; $(1;5,5,3)$, $[0,2]$; $j=26$ & $\tfrac{25}{2}$ & $j\le24$ & level & lifted at 1 loop\\
$\USp(6)$ & none at $j=20,22,24,28$ & -- & $j=20$--$24$, $28$ & -- & --\\
$\Gtwo$ & none for $j\le28$ & -- & $j\le28$ & level $\ge30$ & --\\
\bottomrule
\end{tabularx}
\caption{Lowest-level non-graviton classes in the BMN sector, sector labels $(2J;q_1,q_2,q_3)$ and $\SU(3)_R$ irreps
$[q_1-q_2,q_2-q_3]$. $^a$Given the superconformal selection rule of section~\ref{sec:selection}. $^b$Protected at every
coupling by the index, given the $\SU(2)$ classification. $^c$Order 4 uses the selection rule; order 2 is open
(section~\ref{sec:su3pert}). The $\USp(6)$ scans do not cover $j<20$ or $j=26$; $\USp(6)$ has no class in the sector
$(1;3,3,3)$ at $j=18$ either (\R{05}). For $\SO(9)$ and $\SO(11)$ the lowest energy is not established.}
\label{tab:thresholds}
\end{table}

The statements that differ from the published literature are collected, with equation numbers, in
section~\ref{sec:literature}. The most consequential ones:
\begin{itemize}
\item The criterion ``a multi-graviton polynomial that vanishes on diagonal matrices is $Q$-exact'' (sec.~3 of
\cite{CCKLL}) is false. Counting gravitons on the diagonal undercounts them by the classes that vanish on the Coulomb
branch, and the ``non-graviton index'' built on that count is not a detector of black holes (\R{02}, \R{15}).
\item The $\SU(3)$ threshold operator, eq.~(4.8) of \cite{CCKLL}, is $Q$-exact at $N=3$; the $-t^{24}$ of their
non-graviton index is a fermionic graviton class, represented by their own ``trace relation'' $R_{24}^{(0,3)}$
(\R{09}). The first $\SU(3)$ black hole is at $j=30$ (\R{07}), the brane-fused operator $O^3_{30}$ of \cite{dMK}
(\R{16}); the first class not fused by scalar branes is at $j=34$ (\R{24}).
\item The $\SU(4)$ entries $-\chi_{[2,0]}t^{28}$ and $\chi_{[1,1]}t^{30}$ of eq.~(3.15) of \cite{CCKLL} are artifacts of
the same count (\R{08}, \R{14}).
\item The descendant factor printed in eqs.~(3.7), (3.14), (3.15) of \cite{CCKLL} should be $\prod_I(1-t^2x_I)$
(\R{01}); the thesis \cite{LeeThesis} has the correct factor.
\item The operator $O^3_{30}$ of \cite{dMK} needs three corrections to its printed form to be closed; with them its
stated coefficient $-\tfrac{648}{5}$ is reproduced (\R{16}).
\end{itemize}
Everything else that we checked agrees with the literature, often to the last digit: the $\SU(2)$ operator of
\cite{CKLP}; Table~2 of \cite{CCKLL}; the $\SO(7)$/$\USp(6)$ numbers of \cite{GLRT} and Table~1 of \cite{CLSduality};
the $\SO(7)$ lifting and the rank 741 of \cite{ChoiLee}; the one-loop verdicts and the rank 912 of \cite{KKLP}; and the
entropy excess 1.21\% of \cite{KKLP}.

\section{Setup, definitions and method}\label{sec:method}

\subsection{Objects}\label{sec:setup}

We work with the classical (tree-level) supercharge $Q\equiv Q_0$ of Chang and Lin \cite{CLWords} acting on
gauge-invariant words in the $\frac1{16}$-BPS letters, mostly in the BMN subsector with letters $\ph_m$, $\ps_m$, $\ff$
($m=1,2,3$), on which
\begin{equation}
Q\ph_m=0,\qquad Q\ps_m=\epsilon_{mnp}\,\ph_n\ph_p,\qquad Q\ff=\textstyle\sum_m[\ph_m,\ps_m]
\end{equation}
in the normalisation of the code (\file{code/fortuity/bmn/words.py}; the conventions are those of
\cite{CLWords,CCKLL}, checked against the published counts \R{01}, \R{05}, \R{07}). A BMN sector is labelled
$(2J;q_1,q_2,q_3)$ with $q_I=R_I+J$; the letters carry $\ph_I\to(0;e_I)$, $\ps_I\to(1;(1,1,1)-e_I)$,
$\ff\to(2;(1,1,1))$. The \emph{level} $j=2\sum_Iq_I$ is the power of $t$ in the BMN index (fugacities
$t^jx^{q_1-q_3}y^{q_3-q_2}$), the energy is $E=j/2-J$, and the number of letters is $Y=\sum_Iq_I-2J$. Every letter has
$2J\le j/3$, so $E\ge j/3$ at level $j$. $Q$ commutes with $\SU(3)_R$ and lowers $2J$ by one. An $\SU(3)_R$ highest
weight $\lambda=(q_1\ge q_2\ge q_3)$ labels a \emph{column}, the irrep $[q_1-q_2,q_2-q_3]$ at all $2J$.

Outside the BMN sector we use the full set of letters (gaugini $\lambda$ and derivatives) in the superfield
conventions of \cite{CLWords,BK}, with sectors $(2J_1,2J_2;2q_1,2q_2,2q_3)$ at letter number $Y$; the derivative
grading $d_{\rm BMN}$ vanishes exactly on the BMN sector.

\subsection{Gravitons and fortuity: the operational definition}

A cohomology class is \emph{multi-graviton} (monotone) if it lies in the span of products of single-trace cohomology
generators, modulo $Q$-exact operators, decided by linear algebra with the letters taken as \emph{generic} matrices
of the gauge algebra (not diagonal ones). In the BMN sector of $\SU(N)$ the generators are the components
$u_n,v_n,w_n$ of $\mathrm{tr}\,\Phi^n$; we use $n$ up to (at least) the top Casimir degree of the group ($N$ for
$\SU(N)$, $2N$ for $\SO(2N+1)$ and $\USp(2N)$, 6 for $\Gtwo$), and larger sets where stated. The complement is called
\emph{non-graviton}; it is our operational form of ``fortuitous'' (black-hole candidate) in the sense of
\cite{CLWords,CLCovering}. Consistently, every non-graviton representative that was tested fails to be $Q$-closed at
the next rank: the $\SU(3)$ class of \R{07} at $N=4,5$, the $\SO(7)$ class on $\mathfrak{so}(9)$ (\R{10}), the $\SO(9)$
class on $\mathfrak{so}(11)$ and $\mathfrak{so}(13)$ (\T{43}), the $\SO(11)$ classes on $\mathfrak{so}(13)$ (\R{20}).
(This tests a representative; a different representative closed at higher rank is excluded only by the graviton
test itself.) We do \emph{not} use the
eigenvalue (diagonal-matrix) criterion of \cite{CCKLL}; section~\ref{sec:su3j24} shows why. A class that is a
multi-graviton but vanishes on Cartan-valued letters is called \emph{non-Coulomb} (terminology of \cite{ChoiLee}).

\subsection{The highest-weight complex}

Since $Q$ commutes with $\SU(3)_R$, the complex restricts to highest-weight vectors (the kernel of the raising
operators $E_1,E_2$), and the cohomology of that subcomplex at $(2J;\lambda)$ is the multiplicity of the irrep
$\lambda$ in $H(Q)$ at that $2J$. The dimension $m$ of each highest-weight space is known exactly from an $n_f$-refined
Haar integral and the Weyl character formula. Highest-weight vectors are computed shape by shape in the free
multi-trace space ($E_i$ preserve the flavour-stripped shape of a word). Operators are evaluated on random matrices
of the gauge algebra over $\mathbb F_p$, $p=2^{26}-5$, with fermions expanded in Grassmann frames (one frame per
fermion number), and ranks are taken modulo $p$. The multiplicity of the irrep in cohomology is
$h=m-r_{\rm in}-r_{\rm out}$; its multi-graviton part is the rank of the graviton highest-weight vectors modulo the
exact subspace. Implementation: \file{code/fortuity/hwcoh} (driver \file{run_sectors.py}, parallel engine
\file{sectorpar.py}, FLINT/BLAS kernels \file{ffblas.py}); groups $\SU(N)$, $\SO(2N+1)$, $\USp(2N)$ and $\Gtwo$.

\subsection{Certification}\label{sec:cert}

\begin{enumerate}[label=(\alph*)]
\item \emph{Per-sector self-check.} The rank of the highest-weight vectors at the chosen functionals must equal the
exact multiplicity $m$, separately in every fermion degree. This verifies both that the functionals are faithful on
the highest-weight space and that the word set (often restricted in trace length) spans it. With faithful
functionals the ranks $r$ and the graviton counts are exact over $\mathbb F_p$.
\item \emph{Euler characteristic.} In every column $\sum_{2J}(-1)^{2J}h$ equals the multiplicity in the exact index.
The rank terms cancel in this sum, so it checks the dimensions and the bookkeeping, not the ranks.
\item \emph{Completeness by enumeration.} Every (column, $2J$) with a nonempty letter content is present.
\item \emph{Direction of errors.} A rank computed by random evaluation can only underestimate the true rank
(Schwartz--Zippel \cite{Schwartz,Zippel}). An underestimated rank can create a spurious non-graviton class, never
hide one. Null results (``no non-graviton class'') are therefore robust; positive results were recomputed with
independent seeds and, where stated, with larger generator sets or independent engines.
\item \emph{Rational reconstruction.} Where exact rational coefficients are claimed, they were solved at several
primes and combined by the Chinese remainder theorem \cite{vMS,Peraro}, then re-verified at fresh primes.
\end{enumerate}

\subsection{Other tools}\label{sec:tools}

\begin{itemize}
\item \emph{Direct cohomology} on the full weight space or its Weyl-symmetric part:
\file{code/fortuity/bmn/symdirect.py}, \file{direct.py} (\R{01}, \R{02}, \R{05}, \R{06}).
\item \emph{Exact indices} by Haar integration: \file{code/fortuity/bmn/fullindex.py} ($\SU(N)$),
\file{groupindex.py} ($\SO$, $\USp$, $\Gtwo$), \file{fugacityindex.py} (character values at fixed fugacities).
\item \emph{Diagonal (eigenvalue) graviton counting} of \cite{CCKLL,KMMR}: \file{code/fortuity/bmn/gravindex.py}.
\item \emph{Coulomb test}: rank of the graviton highest-weight vectors on Cartan-valued letters, where $Q$-exact
operators vanish: \file{code/fortuity/hwcoh/noncoulomb.py}.
\item \emph{Block restriction} (partial Coulomb branch): \file{code/fortuity/bmn/blockrestrict.py},
\file{classify.py}, \file{code/fortuity/hwcoh/classreps.py}, \file{classreps_par.py}, \file{cores.py}.
\item \emph{Charge conjugation} $C:X\to-X^T$ and twisted indices: \file{code/fortuity/bmn/cparity.py},
\file{twistedindex.py}.
\item \emph{Full letters and one loop}: \file{code/fortuity/full} (letters $\lambda$ and derivatives, $\SU$, $\SO$,
$\USp$; the one-loop supercharge $Q_1$ of Budzik and Kulp \cite{BK} in \file{q1.py}, with the $\mathfrak{so}(n)$
completeness term; lifting tests in \file{lifting.py}; full-letter levels in \file{sector.py}, \file{levels.py},
\file{gravitons.py}). The parallel and GPU one-loop engine used for the large $\SU(2)$ and $\SU(3)$ tests is in
\file{scripts/oneloop/su-threshold-classes} (\file{fullpar.py}, \file{gpu_rank.py}, \file{cpu_rank.py}).
\item \emph{Two loops}: the two-loop supercharge on derivative-free words, \file{q2bmn.py} in the same directory
(section~\ref{sec:q2}).
\end{itemize}

\section{$\SU(2)$}\label{sec:su2}

\paragraph{The threshold class (\R{01}).} The first non-graviton class of $\SU(2)$ in the BMN sector is unique, at
$j=24$ in the sector $(5;4,4,4)$ ($R_I=\tfrac32$, $J=\tfrac52$, $E=\tfrac{19}{2}$), as found in \cite{CLWords}; we call
it $O_0$. The explicit operator of \cite{CKLP} was verified closed and not exact, and its $SU(1|3)$ descendants
$Q^m_+O_0$ sit at $(6;5,4,4)$ and permutations ($j=26$, bosonic, $[1,0]$) and $(7;5,5,4)$ and permutations ($j=28$,
fermionic, $[0,1]$). The supercharges $Q^m_+$ carry weight $t^2x_m$, so the descendant factor that multiplies the core
of the non-graviton index is $\prod_I(1-t^2x_I)$, and with it the $\SU(2)$ core is exactly $-\sum_{n\ge0}t^{24+12n}$.
Three independent computations agree: direct cohomology with full $2\times2$ matrices in all 683 sectors with $j\le28$;
the exact index (signed sums at every $j\le28$); and the eigenvalue graviton count, which equals the closed form (3.4)
of \cite{CCKLL} through $j=44$.

\paragraph{Complete scan to $j=44$ (\T{27}).} All 446 $\SU(3)_R$ columns with $j\le44$ (3182 sectors, every
nonempty $2J$) were computed with the highest-weight engine on a CPU cluster. Self-checks pass in every sector and the
Euler characteristic equals the exact index in every column. The cohomology has 186 multiplets: 166 multi-graviton,
all of Coulomb type (no non-Coulomb graviton in $\SU(2)$ through $j=44$), and 20 non-graviton classes in 17 sectors:
\begin{center}\small
\begin{tabular}{lcccccccccccc}
\toprule
$j$ & 24 & 26 & 28 & 30 & 32 & 34 & 36 & 38 & 40 & 42 & 44 & total\\
\midrule
non-graviton classes & 1 & 1 & 1 & 1 & 1 & 2 & 3 & 2 & 2 & 3 & 3 & 20\\
\bottomrule
\end{tabular}
\end{center}
All 20 have an odd number of letters, i.e.\ parity $P=-1$ (section~\ref{sec:parity}), and in every column their
number equals $-T$, the kernel parity index computed before the scan finished (\R{17}). The sectors are listed in
\file{data/su2/bmn_hw_cohomology_j44.json} (assembled from the per-sector records by
\file{fortuity.hwcoh.to_json}); the non-graviton counts of all 17 sectors were reproduced by the full-letter one-loop
runs below (different code, different seeds).

\paragraph{One loop (\T{40}, \T{37}).} Each of the 20 classes was tested with the one-loop supercharge $Q_1$ of
\cite{BK}: its image is checked for $Q_0$-exactness in the partner sector $(2J-1,2J-1;2q)$ at $d_{\rm BMN}=1$, with
exact controls, two seeds, and a certificate that the words span. Five classes lift: the multiplet of $O_1$ (one
class each at $j=36$ $[0,0]$, $38$ $[1,0]$, $40$ $[0,1]$, $42$ $[0,0]$) and, at $j=44$ $[1,0]$, one of the two
classes of the sector $(11;8,7,7)$, which has the charges of $w_2O_1$. Fifteen survive: $O_0$ dressed by powers of $w_2$ and
$O_1'$, with their $\mathfrak{psu}(1,2|3)$ descendants ($O_0$'s multiplet at $j=24$--$30$, $w_2O_0$ at $j=32$, $O_1'$ at
$j=36$, $w_2^2O_0$ at $j=40$, $w_2O_1'$ at $j=44$, and the others listed in
\file{scripts/oneloop/su-threshold-classes/runs/t040_oneloop_table.json}). This reproduces
every verdict of \cite{KKLP} in its range ($O_0$, $w_2O_0$, $O_1'$ unlifted; $O_1$ lifted, with the exact-space rank 912
at $j=36$) and extends it to $j=44$. The verdicts had been predicted a priori from the selection rules of
section~\ref{sec:selection} (\T{37}); the prediction was posted before the direct results and agrees in all 16 sectors
it covers. Rank stability was checked at three times the number of evaluation points for the lifted sectors.

\paragraph{All orders for $O_0$.} The order-$n$ partner sectors of $O_0$, $(4,4;8,8,8)$ at $Y=8-n$ with $n\le4$,
have cohomology $0, 27, 0, 4$ for $n=1,2,3,4$, all of Coulomb type. The Coulomb-type classes restrict injectively to
the Cartan, where for $\SU(2)$ the module is the graded-symmetric algebra on one species of letters; that module has
no lowest-weight vector of the weight of $O_0$ at $n=2$ or $4$. So, given the superconformal selection rule of
section~\ref{sec:selection}, $O_0$ has no admissible lifting partner at any perturbative order, with no Coulomb-branch
assumption. Independently, the two-loop differential vanishes by direct computation: $d_2[O_0]=0$ with the $Q_2$ of
section~\ref{sec:q2}, and likewise for the three descendants of $O_0$ at $j=26$, 28, 30
(\file{scripts/oneloop/su-threshold-classes/runs/su2_O0_d2.json} and \file{t040_twoloop_table.json}). The two-loop runs for the other $\SU(2)$ survivors
($j\ge32$) were not completed.

\paragraph{Index protection.} No Coulomb-type multiplet with the weight of $O_0$ and the opposite statistics exists at
any odd rung of its index tower (\file{scripts/oneloop/su-threshold-classes/su2_index_tower.py}, output \file{runs/su2_O0_index_tower.txt}). Given the $\SU(2)$
classification of non-Coulomb classes at these charges (our BMN scans; a direct computation at rung $+1$; for the other
rungs the results of \cite{CLSduality}, not rederived), the index alone protects $O_0$ at every coupling
(section~\ref{sec:selection}).

\paragraph{Higher survivors.} For $w_2O_0$ and $O_1'$ the order-2 and order-4 partner sectors contain only Coulomb-type
classes ($h=119$ and $470$ for $w_2O_0$; $103$ and $1123$ for $O_1'$, with 12 more at order 6) and order 5 is empty. At
order 3 the partner sector of $w_2O_0$ has $h=42{,}894-23{,}094-19{,}797=3$ classes, all non-Coulomb, so its stability
at order 3 is open; the order-3 test of $O_1'$ was not completed.

\section{$\SU(3)$}\label{sec:su3}

\subsection{Level 24: the published threshold is not a black hole}\label{sec:su3j24}

\paragraph{All 29 classes are multi-gravitons (\R{02}).} In the sector $(3;4,4,4)$ ($R=\tfrac52$, $J=\tfrac32$,
$E=\tfrac{21}{2}$) the Weyl-symmetric subspace has dimension 1898, 953 closed and 924 exact operators, hence 29
classes, as in Table~2 of \cite{CCKLL}. The 807 symmetrised graviton monomials have rank 28 on diagonal matrices
(the paper's ``\#gravitons $=28$'') but rank 29 modulo $Q$-exact operators with full $3\times3$ matrices, and
$924+29=953$ saturates the closed space. Every class of the sector is therefore a multi-graviton. Explicit
counterexample to the eigenvalue criterion: a 32-term polynomial $\sum c\,u_2u_2v_2w_3$ with $c\in\{\pm1,\pm2,\pm3\}$,
$Q$-closed, identically zero on diagonal matrices and not exact (\file{data/su3/graviton_class_vanishing_on_diagonal_j24.json};
check: \file{python3 -m fortuity.bmn.r002_check}).

\paragraph{The threshold operator is exact (\R{09}).} The operator $O=-6\,O_6^{(0,3)}$ of eq.~(4.8) of \cite{CCKLL}
was built in two ways (from the printed expansion, and from the definitions through $r_{10},r_{14}$ with the traceless
projections of their (A.3)); the two agree term by term (3747 canonical multi-trace terms). It is closed on
traceless $3\times3$ matrices (not formally, not at $N=4,5$) and it is $Q$-exact at $N=3$, by three independent
computations (the Weyl-symmetric subspace at two primes, the full complex with certified functionals, and the
highest-weight complex). An explicit preimage: an $\ff$-free $\SU(3)$-singlet $\lambda_0$ in $(4;4,4,4)$, a
combination of 99 index-contracted multi-trace monomials with rational coefficients (common denominator 360), with
$O=Q_{\rm code}(\lambda_0)=i\,Q_{\rm paper}(\lambda_0)$ on $3\times3$ traceless matrices
(\file{data/su3/cckll_eq4p8_preimage_j24.json}; test \file{code/tests/test_r009_preimage.py}). The one-dimensional
singlet cohomology of $(j,2J)=(24,3)$ is represented by the paper's own ``fundamental trace relation''
$R_{24}^{(0,3)}=\epsilon_{a_1a_2a_3}u^{a_1b}{v^{a_2}}_b{v^{a_3c}}_d{v^d}_c$ of its (A.4), which vanishes on diagonal
matrices at $N=3$ (not at $N=4,5$) but is not exact; $24R_{24}^{(0,3)}+\mathrm{Sym}(P)$ is exact, with $P$ the
polynomial of \R{02}. Two errors cancel in the paper's Table~2: one graviton too few (28 instead of 29, from the
diagonal count) and one exact operator called non-exact.

\paragraph{The list (A.4) of \cite{CCKLL} (\R{09}).} Of the entries for which the paper gives no $Q$-exact form,
$R_{18}^{(0,2)}$, $(R_{20}^{(0,3)})^i$, $(R_{22}^{(0,3)})^{ij}$ and $(R_{30}^{(0,5)})^{ijk}$ are exact (genuine
relations), while $R_{24}^{(0,3)}$, $(R_{26}^{(0,4)})^i$ and the singlet $R_{30}^{(0,5)}$ are not: they represent the
unique (graviton) classes of their sectors, 5 of the 287 components of (A.4). The entry of level 26 is the descendant
of the one of level 24:
\begin{equation}
(R_{26}^{(0,4)})^1=-\tfrac13\,Q^1_+R_{24}^{(0,3)}+Q(\cdots)\qquad(N=3)
\end{equation}
in the paper's normalisation (\file{scripts/bmn/su3-level24/descendant_check.py}). In the sector of the $[3,0]$ entry
(\R{07}'s sector, four graviton classes and one non-graviton class) the paper's diagonal-vanishing polynomial is
trivial, so the $-t^{30}[3,0]$ of the non-graviton index is the genuine class of \R{07}.

\paragraph{How the three entries vanish on the diagonal.} On diagonal traceless $3\times3$ matrices the three
eigenvalue vectors of the scalars span only a plane in flavour space, and flavour singlets see the fermions through
four anticommuting pairings. As a result the charge sectors of $R_{24}^{(0,3)}$, $(R_{26}^{(0,4)})^i$ and the singlet
$R_{30}^{(0,5)}$ contain \emph{no} Weyl-invariant function of the eigenvalues at all (a Howe-duality count), so every
operator of these sectors vanishes on the diagonal, graviton or not, exact or not; for the $[3,0]$ entry six such
functions exist and the vanishing is a genuine condition (\file{scripts/bmn/su3-level24/diag_vanishing.py}; the counts
are exact linear algebra on polynomials).

\paragraph{Table~1 of \cite{CCKLL} revisited.} With the descendant factor corrected to $\prod_I(1-t^2x_I)$, the
paper's Table~1 (the core of the $\SU(3)$ non-graviton index) is reproduced entry by entry through $j=40$ from our
diagonal counts (\file{table1_core.py}); with the printed factor the rows 26 and 28 would not be empty. Adding the parity
$P$ of section~\ref{sec:parity}: the $[0,0]$ entries at $j=24$ and $30$ are the graviton classes $R_{24}^{(0,3)}$ and
$R_{30}^{(0,5)}$; every other entry through $j=40$ is $P$-odd, hence not a graviton (certified at $j\le30$, and at
$j=32$ by the scan below; conditional on the diagonal counts for $j=34$--$40$); the $[0,0]$ at $j=36$ and the $[1,0]$
at $j=38$ are $P$-even and undetermined. The bookkeeping of the first entry: in the singlet column at $j=24$ the exact
index is $0=(-1)+(+1)$ from the fermionic $R_{24}^{(0,3)}$ at $2J=3$ and a bosonic graviton at $2J=4$; the diagonal
count sees only the latter, so the paper's $Z-Z_{\rm grav}$ has $-1$, which it read as a fermionic black hole in
$[0,0]$. It is the graviton $R_{24}^{(0,3)}$.

\subsection{Level 30: the first black hole}\label{sec:su3j30}

\paragraph{Complete scan (\R{07}).} All $\SU(3)_R$ columns with $j\le30$ (172 columns, every $2J$) were computed: the
first non-graviton class is unique, at $j=30$ in $[3,0]$ (highest weight $(7,4,4)$), $2J=5$, i.e.\
$R=(\tfrac92,\tfrac32,\tfrac32)$, $J=\tfrac52$, $E=\tfrac{25}{2}$, fermionic, in a sector with four graviton classes.
There is no non-graviton class at any $j\le28$. Checks: self-checks in every sector, Euler characteristic $=$ index in
172 of 172 columns, a second seed on the class and its neighbours, graviton generators up to $n=4$ and $n=5$ (\T{21})
with unchanged counts, and agreement with the direct engine on overlapping sectors.

\paragraph{Representative (\T{21}).} An explicit highest-weight representative over $\mathbb F_p$ with 201{,}564
words, closed at $N=3$; $Q$ of its lift is nonzero at $N=4,5$
(\file{scripts/bmn/su3-level30-representative}).

\paragraph{Brane fusion (\R{16}).} Under the block restriction $\SU(3)\supset\SU(2)\times U(1)$ (a chain map for $Q_0$,
section~\ref{sec:fusion}) the class restricts to $x_1^3\otimes O_0$, all other powers of $x_1$ being exact or zero. It
is therefore the brane-fused operator $O^3_{30}$ of \cite{dMK}, eqs.~(2.24)--(2.25), up to normalisation, gravitons and
exact terms. We rebuilt their operator (\file{code/fortuity/bmn/dmk.py}; 6139 terms in its $[3,0]$ highest-weight
component); three readings of the printed formulas are needed, each forced by a check: the mixed-index $r$ and $R$
tensors of their (A.2), (A.3) must be traceless parts; a Koszul sign is needed on the five terms with a fermion-odd
prefactor $vvv$; and $r^{(0,2)}_{16}$ in the $\tfrac85$ term must read $r^{(0,3)}_{16}$. With them the operator is
closed at $N=3$, not at $N=4$, and its $x_1^3$ coefficient is $-\tfrac{648}{5}O^2_{24}$, as stated in their (2.25).

\paragraph{Charge conjugation (\R{17}).} The class is odd under $C:X\to-X^T$, so $P=-1$ (two independent checks: via
the restriction, and by diagonalising $C$ in the sector). In its column the ordinary index is $0$ while the
$C$-twisted index is $-10$.

\paragraph{Lowest energy.} The scans certify the lowest level. The lowest energy $E=\tfrac{25}{2}$ also needs the
sectors with $E<\tfrac{25}{2}$ at higher level, i.e.\ $j=34$ with $2J\ge10$ and $j=36$ with $2J=12$; these small
sectors were not computed (all of $j\le32$ is, see below).

\subsection{Level 32: three classes, all fused}\label{sec:su3j32}

The complete scan of $j=32$ (\T{42}; all 30 columns, 207 sectors, every $2J$, the settings of \R{07}) finds 760
multiplets, 757 multi-graviton and 3 non-graviton: $(6;7,5,4)$ in $[2,1]$ and $(5;8,4,4)$, $(6;8,4,4)$ in $[4,0]$.
Self-checks pass in every sector and the Euler characteristic equals the index in all 30 columns (re-verified when
this repository was assembled). This is exactly the count forced by the kernel parity index of \R{17} ($T=-1$ at
$(7,5,4)$, $T=-2$ at $(8,4,4)$). All three classes are fused (\R{24}\,(iv)): $(5;8,4,4)$ has a non-graviton image
$x^4\otimes O_0$ (the operator $O^3_{32}$ of \cite{dMK}), and $(6;8,4,4)$, $(6;7,5,4)$ have non-graviton images in
their $x^3$ channels. By their charges the last two are the $(1;e_1)$ and $(1;e_2)$ descendants of the class of
\R{07}. Data: \file{data/su3/hw_cohomology_j32_records}.

\subsection{Level 34: the first class beyond scalar brane fusion}\label{sec:su3core}

\paragraph{Counting (\R{24}).} The complete column $\lambda=(8,5,4)$ ($[3,1]$, $j=34$, every $2J$) has 37 multiplets:
34 multi-graviton and 3 non-graviton, two at $(6;8,5,4)$ ($E=14$, 11 letters) and one at $(7;8,5,4)$
($E=\tfrac{27}{2}$). Under $\SU(3)\supset\SU(2)\times U(1)$ the image of a non-graviton class lies in
$x^n\otimes(\text{$\SU(2)$ non-graviton classes})$, equivariantly; the only channels with the right charges are
$x^4\otimes$ (the $\SU(2)$ class at $(6;5,4,4)$) and $x^3\otimes$ (the one at $(7;5,5,4)$), so the fused capacity is
one class at $2J=6$ and one at $2J=7$. Two classes at $2J=6$ against capacity one force at least one class whose
restrictions are all graviton or exact: a class not obtained by attaching scalar branes to an $\SU(2)$ black hole (the
restriction maps of \R{16} keep only the abelian components of the scalars; the abelian components of $\psi$ and $f$
are set to zero).
Independently of the graviton counts, the $C$-twisted index gives exactly $(37-31)/2=3$ $P$-odd classes in the
column, more than the total capacity 2.

\paragraph{Which class it is.} Explicit representatives of all non-graviton classes at $j=30$--$34$ were built
and restricted in every capacity channel (\file{data/su3/core_restrictions_j30_j34.json}): the $2J=7$ class and all
three $j=32$ classes are fused; at $(6;8,5,4)$ the restriction map on the two classes has rank one, so exactly one
combination has no scalar-brane channel. Its images have zero residue in both channels, re-checked on the combined
expression; that class is a generic lift on 834{,}657 of the sector's 834{,}972 trace words, closed at fresh $N=3$
points. With \R{07} fused and the $j=32$ classes fused, \emph{every $\SU(3)$ black-hole class below $j=34$ is
scalar-fused, and the first class that is not sits at $j=34$}, exactly one in this column. By charges, the fused $2J=6$
class sits where the $(1;e_2)$ descendant of $O^3_{32}$ would, and the other one is a new $SU(1|3)$ primary (a reading
by charges, not constructed). This is a statement about scalar branes only. A brane could also carry fermionic
components, and a restriction that keeps the abelian components of $\psi$ would test whether this class is fused in
that wider sense; that extension of the maps of \R{16} has not been done, so whether any $\SU(3)$ black hole is free
of brane structure altogether remains open.

\subsection{Beyond: what the indices force}\label{sec:su3beyond}

The kernel parity index $T$ (section~\ref{sec:parity}) gives lower bounds on the number of $\SU(3)$ black-hole
multiplets: at least $1,3,5,11,19,29$ at $j=30,32,\dots,40$ (and $45$, $66$ at $j=42,44$). These are certified at
$j\le32$ by the direct scans, and checked in the one column $(8,5,4)$ at $j=34$ (exactly 3 $P$-odd classes, graviton
count $34=D$); beyond that they are conditional on the diagonal Coulomb counts not being overcounted (very likely
right, not certified). They include $T=-2$ in the $[n,0]$ column at $j=24+2n$ for $n=4,\dots,8$, the tower of
\cite{dMK}. Independently of any Coulomb count, the $C$-twisted index alone forces at least 4 and 10 $P$-odd (hence
non-graviton) classes in the columns $(9,6,6)$ at $j=42$ and $(10,6,6)$ at $j=44$.

\subsection{Perturbative fate of the first black hole}\label{sec:su3pert}

The class of \R{07}, called R-007 below, sits at $(2J_1,2J_2;2q)=(5,5;14,8,8)$ with $Y=10$ letters,
$d_{\rm BMN}=0$; it is $C$-odd, an $\SU(2)_R$ singlet and a lowest-weight vector of $\mathfrak{psu}(1,2|3)$
(\file{scripts/oneloop/su-threshold-classes/runs/r007_lowest_weight.json}). The loop corrections $Q_n$ have the charges of $Q$, change the letter number by
$1-n$ and raise $d_{\rm BMN}$ by $n$. Hence R-007 is never the target of a lifting, its order-$n$ target lies in
$(4,4;14,8,8)$ at $Y=11-n$, only $n\le4$ can occur, and the target must be $C$-odd.
\begin{itemize}
\item \emph{Order 1} (\T{31}). $Q_1$ of the class is $Q_0$-exact: in the partner sector ($[3,0]$ highest-weight
multiplicity 145{,}197) the exact subspace has rank 76{,}781 at 145{,}475 functionals, and adding $Q_1$ of the class
(or of any of the four gravitons of its sector) does not raise it; controls exact; two independent sets of points; GPU
and CPU arithmetic. R-007 is not lifted at one loop.
\item \emph{Order 3.} The $C$-odd target cohomology is empty, $131{,}048-67{,}869-63{,}179=0$ (exact multiplicity from
twisted Weyl integration; computed ranks are lower bounds, so the zero is rigorous). No lifting at order 3, with no
assumption.
\item \emph{Order 4.} The target has 177 classes, all $C$-odd and all of Coulomb type, and none is a lowest-weight
vector of the weight of R-007; excluded by the selection rule.
\item \emph{Order 2.} The $C$-odd target cohomology has dimension $218{,}008-112{,}730-103{,}379=1{,}899$, of which
$1{,}898$ are of Coulomb type. The one non-Coulomb class is not a lowest-weight vector (its $[2,1]$ and $[2,0]$ lowering
images are not exact, with span certificates), so the selection rule excludes it. Eight Coulomb-type classes pass all
lowest-weight tests; $d_1$ vanishes on them. They are excluded under the Coulomb-branch assumption (a Coulomb-type
class is never a target), and otherwise only by $d_2[\text{R-007}]$. The direct computation of $d_2[\text{R-007}]$
with the $Q_2$ of section~\ref{sec:q2} was set up on the cluster in two independent chains; it has not produced an
exactness verdict, so order 2 is open.
\item \emph{Restriction argument} (\T{33}). Theorem A: every possible order-3 partner of R-007 is an $\SU(3)$ core
(restricts to zero under $\SU(3)\supset\SU(2)\times U(1)$); classical, with all fused channels computed directly.
The block restriction is not a $Q_1$ chain map (W-boson exchange), so it can prove lifting but not survival.
\end{itemize}
So R-007 is stable to all orders under the Coulomb-branch assumption, and without it at orders 1, 3 and 4. The index
cannot do better (section~\ref{sec:selection}): the eight multiplets above have the weight and $C$-parity of R-007 and
the opposite statistics, so its protection is equivalent to the Coulomb-branch statement for those eight classes.
The all-letter index at the weight of R-007 is $35$, its $C$-twisted version $-565$; without the all-letter graviton
subtraction these decide nothing.

\paragraph{The non-Coulomb graviton $R_{24}^{(0,3)}$ at loop level.} It can never be a target (grading), and it is
not a source at one loop ($Q_1R_{24}^{(0,3)}$ is $Q_0$-exact, exact-space rank 1897, two sets of points; \T{30}) nor
at two loops (direct null test at rank 3 with the two-loop supercharge). The same pipeline reproduces the one-loop
lifting of one of the two $\SU(2)$ classes at $j=36$ \cite{KKLP}.

\section{$\SU(4)$ and $\SU(5)$}\label{sec:su4}

\begin{itemize}
\item \R{03}. The diagonal (eigenvalue) count reproduces eq.~(3.15) of \cite{CCKLL}:
$Z-Z_{\rm grav}=-\chi_{[2,0]}t^{28}+\chi_{[1,1]}t^{30}+O(t^{32})$. For $\SU(5)$ the same difference vanishes through
$j=28$. No direct $\SU(5)$ scan was done.
\item \R{08}, \R{14}. Complete highest-weight scans of all columns with $j\le28$ (145 columns) and $j=30$ (27
columns), every $2J$: all cohomology is multi-graviton, so the $\SU(4)$ BMN threshold is at $j\ge32$. Both terms of
eq.~(3.15) are artifacts of the diagonal count: the column $(6,4,4)$ ($[2,0]$, $j=28$) has $1,2,6,30,46,19,1$ classes
at $2J=0,\dots,6$ and the column $(6,5,4)$ ($[1,1]$, $j=30$) has $0,3,13,36,66,50,13,1$ at $2J=0,\dots,7$, all
multi-graviton. Self-checks in every sector; Euler characteristic $=$ index in 145/145 and 27/27 columns.
\item Non-Coulomb gravitons (\R{15}): $(j,\lambda,2J)=(28,(6,4,4),3)$, $(30,(7,4,4),4)$, $(30,(6,5,4),4)$,
$(30,(7,4,4),3)$, one class each. The onset $j=28$ is the ``apparent threshold'' of the diagonal index.
\item Parity (\R{17}). Every $\SU(4)$ class through $j=30$ is $P$-even. The $C$-twisted index, validated against the
direct cohomology in all 172 columns with $j\le30$, is non-negative in every column through $j=44$, so the indices
force no $\SU(4)$ black hole in that range (\file{scripts/index/c-parity/runs/su4_twisted_scan_j44.json}).
\item A fused candidate for the first $\SU(4)$ class ($x^3y^3\otimes O_0$ in $[6,0]$ at $j=36$, $2J=5$, by analogy
with the $\SO$ tower) was proposed but not computed; it would bound the threshold level from above only.
\end{itemize}

\section{$\SO(2N+1)$, $\USp(2N)$, $\Gtwo$}\label{sec:so}

\subsection{Index differences}

$\SO(2N+1)$ and $\USp(2N)$ are exchanged by S-duality, and their BMN indices agree at low level. The first
differences (\R{04}, \R{21}; exact Haar integrals, \file{code/fortuity/bmn/groupindex.py}, \file{fugacityindex.py}):
\begin{align*}
I_{\SO(7)}-I_{\USp(6)}&=-t^{18}+3t^{20}-3t^{22}+11t^{24}-33t^{26}+39t^{28}-74t^{30}+\dots\\
I_{\SO(9)}-I_{\USp(8)}&=-6t^{22}+18t^{24}-33t^{26}+117t^{28}-289t^{30}+564t^{32}+\dots\\
I_{\SO(11)}-I_{\USp(10)}&=-21t^{26}+63t^{28}-154t^{30}+\dots\\
I_{\SO(13)}-I_{\USp(12)}&=-56t^{30}+168t^{32}+\dots
\end{align*}
with $\SU(3)_R$ content of the first term $[0,0]$, $[2,0]$, $[4,0]+[0,2]$ and $[6,0]+[2,2]+[0,0]$ respectively, i.e.\
$-\mathrm{Sym}^{N-3}[2,0]$ (dimension $\binom{N+2}{5}$) at $j=4N+6$, the charge of one $\ps$ and $2N+1$ scalars. For
$N=3$ this reproduces \cite{GLRT}; $N=4,5,6$ are new as far as we know. For $N=6$ the content is established by
character values at three random fugacity points plus the unrefined count (strong evidence, not a proof).

\subsection{Direct cohomology at the first difference}

\begin{center}\small
\begin{tabular}{llcccccc}
\toprule
group & sector & dim & closed & exact & classes & multi-graviton & non-graviton\\
\midrule
$\SO(7)$ & $(1;3,3,3)$ & 903 & 683 & 559 & 124 & 123 & 1\\
$\USp(6)$ & $(1;3,3,3)$ & 903 & 682 & 559 & 123 & 123 & 0\\
$\SO(9)$ & $(1;3,3,3)$ & 936 & 699 & 559 & 140 & 140 & 0\\
$\USp(8)$ & $(1;3,3,3)$ & 936 & 699 & 559 & 140 & 140 & 0\\
\midrule
$\SO(9)$ & $(1;5,3,3)$ & 5395 & 4133 & 3704 & 429 & 428 & 1\\
$\USp(8)$ & $(1;5,3,3)$ & 5395 & 4132 & 3704 & 428 & 428 & 0\\
\bottomrule
\end{tabular}
\end{center}
(\R{05}, \R{06}; weight spaces, two seeds each.) The $\SO(7)$/$\USp(6)$ rows reproduce Table~1 (BMN column) of
\cite{CLSduality} and sec.~4.3 of \cite{GLRT}; the $\SO(9)$/$\USp(8)$ rows at 8 letters are a control. The $\SO(9)$
class ($\ps\ph^9$, $j=22$, $E=\tfrac{21}{2}$, highest weight of a $[2,0]$) is new as far as we know; the complete
$\SO(9)$ scan of level 22 (16 columns, 79 sectors: 320 multiplets, 319 multi-graviton) reproduces it as the only
non-graviton class at that level. For $\SO(11)$ (\R{11}, highest-weight columns at $j=26$, two seeds):
\begin{center}\small
\begin{tabular}{lccc}
\toprule
column & $2J$ & $\SO(11)$: classes / graviton / non-graviton & $\USp(10)$\\
\midrule
$(7,3,3)$, $[4,0]$ & 0, 1, 2 & 0/0/0, \ 23/22/1, \ 0/0/0 & 0/0/0, \ 22/22/0, \ 0/0/0\\
$(5,5,3)$, $[0,2]$ & 0, 1, 2 & 0/0/0, \ 12/11/1, \ 0/0/0 & 0/0/0, \ 11/11/0, \ 0/0/0\\
\bottomrule
\end{tabular}
\end{center}
So for $N=3,4,5$ the first index difference is exactly one fermionic non-graviton $\SO(2N+1)$ multiplet per irrep at
$2J=1$, absent in $\USp(2N)$, with all graviton classes shared. The representatives fail to close one rank up: the
$\SO(7)$ class (eq.~(4.14) of \cite{GLRT}) on $\mathfrak{so}(9)$; the $\SO(9)$ representative on $\mathfrak{so}(11)$
and $\mathfrak{so}(13)$ (\T{43}; \file{scripts/bmn/so-usp-g2-scans/extensions/closure_check.py}); the $\SO(11)$ ones on
$\mathfrak{so}(13)$.

\subsection{These are the lowest-level classes (\R{18})}

Complete scans of every column and every $2J$ below the first difference:
\begin{center}\small
\begin{tabular}{lccccc}
\toprule
group & scanned & columns & sectors & classes & non-graviton\\
\midrule
$\SO(7)$ & $j\le18$ & 52 & 170 & 166 & 1 (\R{05}'s class)\\
$\SO(9)$ & $j\le22$ & 82 & 313 & 665 & 1 (\R{06}'s class)\\
$\SO(11)$ & $j\le24$ & 101 & 416 & 1331 & 0\\
\bottomrule
\end{tabular}
\end{center}
With $\SO(3)=\SU(2)$ ($j=24$) and $\SO(5)$ (none for $j\le28$), the lowest fortuitous level of $\SO(2N+1)$ in the BMN
sector is $j=24$, $\ge30$, $18$, $22$, $26$ for $N=1,\dots,5$. For $\SO(7)$ the class is also the lowest in energy
(all levels $j\le28>3\cdot\tfrac{17}{2}$ are scanned); for $\SO(9)$ and $\SO(11)$ the energy statement is not
established. Checks: self-checks in all sectors; Euler characteristic $=$ exact refined index in every column; the
$\SO(9)$ column $(5,3,3)$ at $j=22$ as a positive control (sector by sector $m=131/442/664/665/404/114/3$, one
non-graviton class at $2J=1$); the audit script \file{scripts/bmn/so-odd-thresholds/analyze.py} reproduces the table
from the data files.

\subsection{$\SO(7)$, $\USp(6)$ and $\Gtwo$ to level 28; $\SO(5)$}

\begin{itemize}
\item \R{13}, \R{19}. Complete scans of $\SO(7)$ and $\Gtwo$ for $j=20,\dots,28$ and of $\USp(6)$ for $j=20,22,24$ and
$28$ ($j=26$ was not scanned for $\USp(6)$).
$\SO(7)$ non-graviton classes: at $j=20,22,24$ the $SU(1|3)$ descendants of the $j=18$ class
($(2;4,3,3)$, $(3;4,4,3)$, $(4;4,4,4)$); at $j=26$ a new multiplet $(3;5,4,4)$, $[1,0]$, $E=\tfrac{23}{2}$, which is
the product $O_{18}w_2$ of the $j=18$ class with the graviton $w_2$ (a ``hairy'' class: the sector has exactly one
non-graviton class and $O_{18}w_2$ is a nonzero non-graviton class there); at $j=28$ two classes at $2J=4$ in
$(5,5,4)$ and $(6,4,4)$, with the charges of the descendants of the $j=26$ class. Level totals at $j=28$: $\SO(7)$ 1306
multiplets (2 non-graviton), $\USp(6)$ 1302, $\Gtwo$ 212 (all multi-graviton).
\item $\USp(6)$ has no non-graviton class at $j=20,22,24,28$, nor in the sector $(1;3,3,3)$ at $j=18$ (\R{05});
$j<20$ and $j=26$ were not scanned.
\item $\Gtwo$ (\R{10}, \R{22}): multi-trace words span all $\Gtwo$ invariants in every BMN sector with $j\le18$
(exact Haar dimension $=$ rank of trace words); $\Gtwo$ has no non-graviton class for $j\le28$ (complete scans; 83
classes for $j\le18$). The $\SO(7)$ class restricted to $\mathfrak g_2$-valued letters is a nonzero, non-exact
\emph{multi-graviton} class of $\Gtwo$, and so is every $\SO(7)$ non-graviton class at $j=22$--$26$. Caveat: the
trace-word basis is proven complete for $\Gtwo$ only for $j\le18$.
\item $\SO(5)=\USp(4)$ (\R{12}): complete scan of $j\le28$ (676 sectors, 960 classes), all multi-graviton, so the
threshold level is $\ge30$ and any fortuitous class has $E\ge10$. Four columns were recomputed with the $\USp(4)$
letters, identically.
\item One $\SO(7)$ multi-graviton class without a $\USp(6)$ counterpart, at $(2;6,3,3)$, $j=24$, accounts for the
$+[3,0]t^{24}$ of the index difference; it vanishes on the Coulomb branch (\R{13}\,(v)).
\end{itemize}

\subsection{Brane fusion of the $\SO$ classes (\R{16})}

Under $\SO(9)\supset\SO(7)\times\SO(2)$ the $\SO(9)$ class restricts to $x_1^2\otimes O_{\SO(7)}$ plus gravitons;
under $\SO(11)\supset\SO(7)\times\SO(2)^2$ (blocks $x,y$) the two $\SO(11)$ classes restrict to
$x_1^2y_1^2\otimes O_{\SO(7)}$ ($[4,0]$) and $(x_1y_2-x_2y_1)^2\otimes O_{\SO(7)}$ ($[0,2]$; measured coefficient ratio
$1:-2:1$). The first $\SO(2N+1)$--$\USp(2N)$ difference is the $\SO(7)$ black hole with $N-3$ separated branes, each
carrying one scalar bilinear ($\Delta j=4$, $\Delta E=2$), which accounts for $j=4N+6$, $E=2N+\tfrac52$ and the content
$\mathrm{Sym}^{N-3}[2,0]$; it also explains the absence of the tower in $\USp(2N)$, which has no class at $j=18$. This
gives certificates of fortuity for the $\SO(9)$ and $\SO(11)$ classes that use only the $\SO(7)$ data.

\subsection{One loop (\R{20})}

With the $\mathfrak{so}(n)$ version of $Q_1$ (\file{code/fortuity/full/q1.py}; checked against
$Q_1\,\mathrm{tr}(\ph_1\ps_1)=4(n-2)\,\mathrm{tr}(\lambda_+\lambda_-)$ of \cite{ChoiLee} and the $\mathfrak{so}(n)$
completeness relation): the $\SO(7)$ class is lifted, $Q_1O_{18}$ not lying in the span of the exact space (rank
741, the number of \cite{ChoiLee}) and $Q_1$ of the 324 graviton polynomials of its sector; the $\SO(9)$ class and both
$\SO(11)$ classes are lifted as well, through their $\SO(7)$ core (one-way block tests with 1364, 4436 and 6458
graviton polynomials). In every sector tested, $Q_1$ of every graviton polynomial is $Q_0$-exact. So every
$\SO(2N+1)$ class at the first $\SO$/$\USp$ difference, $N=3,4,5$, is lifted at one loop, as S-duality requires
(\file{python -m fortuity.full.lifting so7|so9|so11}; outputs in \file{scripts/oneloop/so-odd-classes/runs}). The block
restriction commutes with $Q_1$ in cohomology on all three lifted fused classes (\T{33}; ratios $-2$, $288$, $-60$ on
both sides; \file{scripts/oneloop/protection}).

\subsection{Full letters: $\SO(7)$ and $\USp(6)$ at levels 18--22 (\R{25})}

Complete highest-weight cohomology with all letters (gaugini and derivatives) at levels 18, 20, 22 (every column,
every $(2J_1,2J_2)$, every letter number):
\begin{itemize}
\item Level 18: $\SO(7)$ 985 multiplets, $\USp(6)$ 983. They differ only in the two sectors of Table~1 of
\cite{CLSduality}, both reproduced in full: $(1,1;6,6,6)$, $Y=8$ (the BMN class: $903/559$ exact$/124$ classes,
$123+1$, against $903/559/123$) and $(0,0;6,6,6)$, $Y=8$ (the chiral ring: $826/741/85$ against $825/741/84$, all
multi-graviton; the extra $\SO(7)$ graviton vanishes on the Coulomb branch, as stated in \cite{CLSduality}).
\item Level 20: $\SO(7)$ 2588, $\USp(6)$ 2578; ten extra $\SO(7)$ multiplets in nine sectors, all at $Y=8$: five are
non-graviton, five are $\mathfrak{so}(7)$-only multi-gravitons, all Coulomb-invisible.
\item Level 22: $\SO(7)$ 6747, $\USp(6)$ 6713; 34 extra $\SO(7)$ multiplets in 19 sectors, all at $Y=8$; of the 21 in
the twelve sectors that were classified, five are non-graviton and sixteen Coulomb-invisible multi-gravitons; 13 in the
seven largest sectors are not classified.
\item In every classified sector the Coulomb-type graviton counts of $\SO(7)$ and $\USp(6)$ agree, as S-duality
requires. The derivative-free sectors agree column by column with the BMN engine.
\end{itemize}

\subsection{Kernel parity for $\SO$, $\USp$, $\Gtwo$ (\T{34})}

Here $C$ acts trivially on invariants, so $P=(-1)^Y$. The kernel parity index (twisted index minus Coulomb count) for
$\SO(7)$ to $j=30$ and $\USp(6)$, $\Gtwo$ to $j=28$ is never negative: no black hole is forced by indices in these
groups (their known classes have $P=+1$, which the index cannot see). Positive values are non-Coulomb gravitons where
checked directly ($(8,3,3)$ at $j=28$ in $\SO(7)$ and $\USp(6)$; two $\Gtwo$ sectors); elsewhere they rest on the
signed Coulomb count (not certified).

\section{Structural results}\label{sec:structure}

\subsection{Charge-conjugation parity and index-only detection of black holes}\label{sec:parity}

\paragraph{Parity lemma (\R{17}).} Charge conjugation $C:X\to-X^T$ on every letter commutes with $Q$. With $Y$ the
number of letters, every multi-graviton has $P:=C(-1)^Y=+1$: single-trace gravitons are components of
$\mathrm{tr}\,\Phi^n$, which have $n$ letters and satisfy $\mathrm{tr}((-\Phi^T)^n)=(-1)^n\mathrm{tr}\,\Phi^n$, and both
$C$ and $(-1)^Y$ are multiplicative. Hence a class with $P=-1$ is automatically non-graviton. For groups on whose
invariants $C$ is inner or trivial ($\SU(2)$, $\SO(2N+1)$, $\USp(2N)$, $\Gtwo$) the lemma says that every graviton has
an even number of letters; it holds for all 4{,}475 graviton classes of these groups in our data. On Cartan-valued
letters $C$ acts as $(-1)^Y$, so every class visible on the Coulomb branch has $P=+1$.

\paragraph{Twisted index.} $I_C=\mathrm{Tr}\,(-1)^{2J}C$ at fixed $q$ equals $(-1)^{|q|}(\dim H^{P=+1}-\dim H^{P=-1})$,
summed over $2J$ with no boson--fermion cancellation. Computed by twisted Weyl integration over the coset
$\SU(N)\cdot C$ (\file{code/fortuity/bmn/twistedindex.py}), it agrees with the cohomology data in all 172 $\SU(3)$
and all 172 $\SU(4)$ columns with $j\le30$. In the column of R-007 the ordinary index is $0$ and $I_C=-10$.

\paragraph{Kernel parity index.} With $D$ the Coulomb-branch (diagonal) graviton count of a column,
$T=(-1)^{|\lambda|}I_C-D$ is $\#\{P=+1\}-\#\{P=-1\}$ among the classes invisible on the Coulomb branch, so $T<0$ forces
black holes. The independent review (section~\ref{sec:checks}) showed that $D$, a six-term signed Weyl alternation of
randomly sampled dimensions, is not a certified lower bound; with that correction: the $\SU(2)$ values through $j=44$
are exact (in the 17 forced columns $D=0$, and every $\SU(2)$ diagonal dimension was also computed exactly over
$\mathbb Q$); the $\SU(3)$ values are confirmed by direct cohomology at $j\le32$ and in the column $(8,5,4)$; beyond
that they are very likely right but conditional. $Q_1$ preserves $Y$ and commutes with $C$, so it conserves $P$: at one
loop a $P$-odd class can pair only with another $P$-odd (non-graviton) class. This is why the $P$-odd threshold
classes of $\SU(2)$ and $\SU(3)$ survive one loop, while the $P$-even $\SO(7)$ class is lifted by a graviton
\cite{ChoiLee}.

\subsection{Non-Coulomb gravitons}\label{sec:noncoulomb}

Multi-graviton classes that vanish on Cartan-valued letters (\R{15}, \R{23}; complete lists in
\file{data/noncoulomb}). Onsets in the BMN sector: $\SU(2)$ none through $j=44$; $\SU(3)$ at $j=24$ (the multiplet of
$R_{24}^{(0,3)}$, verified explicitly: $R_{24}^{(0,3)}$ and $Q^1_+R_{24}$, $Q^1_+Q^2_+R_{24}$, $Q^1_+Q^2_+Q^3_+R_{24}$
are non-exact graviton classes at $(26,[1,0],4)$, $(28,[0,1],5)$, $(30,[0,0],6)$, plus a second primary, the singlet
$R_{30}^{(0,5)}$ at $(30,[0,0],5)$); $\SU(4)$ at $j=28$; $\SO(5)$ at $j=20$; $\SO(7)$ exactly at $j=24$; $\USp(6)$ at
$j=28$ (none at $j=20$--$24$); none in $\SO(9)$ through $j=22$ or $\SO(11)$ through $j=24$. None occurs at $2J=0,1$ in
the range computed. Properties:
\begin{itemize}
\item \emph{S-duality.} The Coulomb-type graviton counts of $\SO(2N+1)$ and $\USp(2N)$ agree in every sector compared
(99 common sectors for $j\le24$, all of $j=28$, and the classified full-letter sectors at levels 18--22), as their
isomorphic Cartan data require; graviton mismatches between the two groups sit only in non-Coulomb classes.
\item \emph{Rank dependence.} $R_{24}^{(0,3)}$ vanishes on diagonal $3\times3$ matrices but not on $4\times4$ or
$5\times5$ ones; the observation that Coulomb-invisibility is rank-specific, and the terminology, are due to
\cite{ChoiLee}; the first non-Coulomb graviton is the $\SO(7)$ chiral-ring class of \cite{CLSduality}.
\item \emph{Counting.} The finite-$N$ graviton count by ``quantizing the moduli space'' (eigenvalues of commuting
diagonal letters, \cite{KMMR} sec.~6.1, used for the $\frac1{16}$-BPS sector in \cite{CCKLL} sec.~3) undercounts the
monotone spectrum by exactly these classes; $Z-Z_{\rm grav}$ computed that way is the index of the kernel of the
Coulomb restriction, which contains non-Coulomb gravitons as well as black holes.
\item \emph{Quantum.} BMN classes are never targets of a lifting (section~\ref{sec:selection}); $Q_1$ of every graviton
polynomial tested is $Q_0$-exact, in all groups (\R{20}, \T{30}, \T{31}).
\end{itemize}

\subsection{Brane fusion and cores}\label{sec:fusion}

Evaluate an operator on block-diagonal letters in which some abelian blocks carry only scalar components $x$
(separated branes). Since $Q$ is built from commutators and the abelian parts are central, the coefficient of each
monomial $x^a$ defines a chain map from the complex of $G$ to that of the unbroken non-abelian factor; it sends exact
to exact and multi-gravitons to multi-gravitons. A non-graviton image therefore proves that a class is non-graviton and
identifies the smaller-group class inside it (\R{16}). Results: the first $\SU(3)$ class is $x_1^3\otimes O_0$, the
first $\SO(9)$ and $\SO(11)$ classes are the $\SO(7)$ class dressed by one scalar bilinear per separated $\SO(2)$
block, and the first $\SU(3)$ class with no fused image is at $j=34$ (\R{24}). In every case examined the fused classes
share the one-loop fate of their core (the $\SO$ classes are lifted with the $\SO(7)$ core; the $\SU(3)$ class
survives like $O_0$). The restriction is not a chain map for $Q_1$ (W-boson exchange), so it can prove lifting but not
survival; the intertwining in cohomology was checked on three lifted classes (\T{33}). Reading: an $\SO(2)$ block is a
D3-brane with its orientifold image; the $\SU(3)$ class is the $\SU(2)$ black hole bound to one brane, the
``dual-giant hair'' of \cite{dMK}. Thresholds read off as ``first fortuitous class at rank $N$'' therefore mix cores
and fused classes.

\subsection{Selection rules, and what the index protects}\label{sec:selection}

\paragraph{Gradings.} $Q_n$ changes the letter number by $1-n$ and raises $d_{\rm BMN}$ by $n$. So a BMN class is never
the target of a lifting, the order-$n$ partner of a class with $Y$ letters has $Y+1-n$ letters, and only finitely many
orders can occur for a given class. $C$ is an exact symmetry of $\SU(N)$ $\cN=4$ SYM and commutes with $Q(g)$: partners
have the same $C$, and the order-$n$ partner of a $P$-odd class has $P=(-1)^n$, so at odd $n$ it cannot be a graviton.

\paragraph{Superconformal selection rule.} The loop differentials commute with the tree-level generators of the
commutant $\mathfrak{psu}(1,2|3)$ of $Q$; a lifting partner must then be a lowest-weight vector of the same weight. This
follows from perturbative superconformal invariance (a framework hypothesis, not the Coulomb-branch assumption). $Q_1$
commutes exactly with all these generators except $K$ and $T$; for those, the defect was checked to be $Q_0$-exact on
the one-loop partner of the lifted $\SU(2)$ class at $j=36$. At two loops the rule predicts $d_2[O_0]=0$, which the
direct computation confirms. It is argued and tested, not proven, beyond one loop.

\paragraph{The index.} The BPS states at coupling $g$ form a unitary lowest-weight module of $\mathfrak{psu}(1,2|3)$,
whose Cartan charges are exactly the charges kept by the index. Since characters of distinct lowest-weight modules are
independent, for each lowest weight $\mu$ the number $n_\mu=\sum_k(-1)^ka_k(\mu)$, with $a_k$ the number of BPS
multiplets of lowest weight $\mu$ at rung $k$ (sectors $(2J-k,2J-k;2q)$), is independent of the coupling (for
$\SU(N\ge3)$ separately for $C=\pm1$). Hence: at least $|n_\mu|$ multiplets of weight $\mu$ are BPS at every coupling;
multiplets leave the BPS spectrum in pairs of the same $\mu$ and $C$ and opposite statistics; and if no multiplet of
weight $\mu$ exists at the rungs of one parity, all those at the other parity are BPS at every coupling. The last
statement protects the $\SU(2)$ threshold class non-perturbatively (section~\ref{sec:su2}) and fails for the $\SU(3)$
class, for which eight Coulomb-type multiplets of its weight and $C$-parity exist at rung $+1$
(section~\ref{sec:su3pert}).

\subsection{The two-loop supercharge on derivative-free words}\label{sec:q2}

$Q_1^2\neq0$ on words (on 26 of 60 derivative-free test words), while nilpotency of $Q(g)$ requires
$Q_1^2+\{Q_0,Q_2\}=0$ (eq.~(12) of \cite{BK}), so a two-loop term is needed; \cite{BK} compute $Q_1$ only. We derive it
on derivative-free words. In the superfield Feynman rules of the holomorphic twist the $l$-loop term of $Q$ is a sum over
Laman graphs with $l+2$ vertices, one of them the operator \cite{BK}. At two loops the only graph is $K_4$ minus an
edge; with the operator at a two-valent vertex the three interaction vertices form a triangle and give
$\delta^3(0)=0$, which leaves one graph, with three letters in and two out ($\Delta Y=-1$):
\begin{equation}
Q_2\big(C^A(\theta_1)C^B(\theta_2)C^C(\theta_3)\big)=K_2\,f_{APX}f_{BPQ}f_{CQY}\,
\delta^3(\theta_1-\theta_2)\,\delta^3(\theta_2-\theta_3)\,D_2\big(C^X(\theta_1),C^Y(\theta_3)\big),
\end{equation}
summed over the choice of the middle letter. On underived letters the kinematic factor must be an $\SU(2)_R$ invariant
with four derivatives, and there is exactly one,
$D_2(f,g)=\epsilon^{ac}\epsilon^{bd}\,\partial_a\partial_bf\,\partial_c\partial_dg$. So on derivative-free words
$Q_2=K_2U$ with $U$ known, and nilpotency fixes the constant:
\begin{itemize}
\item $Q_1^2+K_2\{Q_0,U\}=0$ holds on 120 random derivative-free multi-traces (one or two traces of 2--5 letters) for
$N=2,3,4$ with the single value $K_2=-\tfrac16$ ($Q_1$ normalised as in \file{fortuity.full.q1}); $Q_1^2\ne0$ on 54
of them, and $\{Q_0,U\}\ne0$, so the value is unique;
\item negative control: changing the weight of one term of $D_2$, no constant works on any of the 54 words;
\item $U$ commutes with $\SU(3)_R$ and $C$ and anticommutes with $\partial/\partial\theta_a$ and $S_a$, as $Q_1$ does,
and maps trace relations to relations at $N=2$ and $N=3$;
\item eq.~(2.19) of \cite{KKLP}, $Q_1(Q_1O_0)=-Q_0(Q_2O_0)$, holds at rank 2 with this $Q_2$ and fails without it or
with another constant.
\end{itemize}
Assumed: the true $Q_2$, in the scheme in which $Q_1$ is the operator above, has the graph form; within that form the
solution is unique, and conjugations of $Q(g)$ change $d_2$ only by $d_1$-exact classes. $Q_2$ on letters with
derivatives is not determined (it is not needed for $d_2$ of a derivative-free class). As far as we know this two-loop
supercharge is not in the literature. Code: \file{scripts/oneloop/su-threshold-classes/q2bmn.py}, \file{test_q2.py},
\file{test_q2_tracerel.py}, \file{d2_make.py}.

\subsection{Entropy-level lifting}\label{sec:entropy}

\begin{itemize}
\item \T{38}. The bonus-charge refinement of the index in \cite{KKLP} implies that a fraction of the classical entropy
must lift; at equal charges we reproduce their excess, $\delta_0=0.012134434$ (1.21\%). Refining by $P$ does not
sharpen the bound: at fixed charges, $J_R$ and $d_{\rm BMN}$ the letter-number parity is fixed, so $P$ reduces to
$\pm C(-1)^F$ and the $P$-graded index is the $C$-twisted one, whose Cardy growth cannot exceed the untwisted one (half
of it for $\SU(N\ge3)$, the same for groups on which $C$ is trivial; for $\SO(7)$ we recover their ratio 1.0121).
\item \T{41}. The charge dependence: over 927 points of the ordered wedge of charge space, $\delta_0$ is largest at
equal charges and decreases along every sampled ray; the supersymmetric black holes fill the disk
$\sum_I\varphi_I^2<1$ in charge fractions, not just the triangle $\varphi_I>0$; along the edge $Q_3=0$, $\delta_0$
falls from $0.0100$ at the midpoint to $0.0011$ at fractions $(0.95,0.05,0)$. To first order in
$\epsilon=(\dim G/j)^{1/3}$ at equal charges, $\delta=0.0121+0.0551\,\epsilon+o(\epsilon)$, and the first-order
coefficient is positive at every computed point. An independent code reproduces $\delta_0$ to $5\times10^{-10}$ and the
first-order coefficient to $5\times10^{-8}$ at 173 common points. The collapsed holonomy saddle dominates in all 432
verdicts (six groups, 18 charge points, four criteria) inside the triangle; dominance was not tested where a charge is
negative. $\delta_0$ bounds the lifted fraction only under the stated saddle-point assumptions.
\end{itemize}

\section{Comparison with the literature}\label{sec:literature}

Paper by paper, every statement of ours that overlaps with a published one. ``Agrees'' means reproduced by our own
computation, to the stated precision.

\begin{description}[style=nextline,leftmargin=1.2em]
\item[Chang, Lin \cite{CLWords}; \cite{CLCovering}.] Conventions, supercharge and the notion of fortuity. Agrees: the
$\SU(2)$ BMN threshold at $j=24$, $(5;4,4,4)$, $E=\tfrac{19}{2}$, unique, with nothing below it in the BMN sector
(\R{01}). The non-graviton representatives we tested fail to close at the next rank, as for fortuitous classes
(section~\ref{sec:method}).
\item[Choi, Kim, Lee, Park \cite{CKLP}.] Agrees: the explicit $\SU(2)$ threshold operator is closed and not exact; its
$SU(1|3)$ descendants are non-exact non-gravitons (\R{01}).
\item[Choi, Choi, Kim, Lee, Lee \cite{CCKLL}.] The overlap is large; item by item:
\begin{itemize}
\item Table~2 (the $j=24$ singlet sector: 1898 / 953 / 924 / 29): agrees exactly. Its split ``28 gravitons + 1 black
hole'' does not: the 29 classes are all multi-gravitons (\R{02}).
\item Sec.~3 and its appendix A, criterion ``graviton polynomial vanishing on diagonal matrices $\Rightarrow$
$Q$-exact'': false; explicit counterexample (\R{02}); in the three affected sectors of (A.4) there is no function on the
diagonal at all (section~\ref{sec:su3j24}).
\item Eq.~(4.8), the threshold representative: $Q$-exact at $N=3$, explicit preimage (\R{09}). Their transcription and
definition of the operator agree with each other (3747 terms); why their test reported it non-exact is not known.
\item Eq.~(3.4), the $\SU(2)$ graviton closed form: agrees through $j=44$.
\item Eqs.~(3.7), (3.14), (3.15): the descendant factor $\prod_{I<J}(1-t^4x_Ix_J)$ is a misprint for
$\prod_I(1-t^2x_I)$ (the thesis \cite{LeeThesis}, eq.~6.33, has the correct one); with it the $\SU(2)$ core is exactly
$-\sum_nt^{24+12n}$ and Table~1 is reproduced.
\item Eq.~(3.15), $\SU(4)$: the diagonal-count terms $-\chi_{[2,0]}t^{28}$, $\chi_{[1,1]}t^{30}$ are reproduced as
diagonal-count statements and are artifacts: no non-graviton class in $\SU(4)$ for $j\le30$ (\R{03}, \R{08},
\R{14}).
\item List (A.4): $R_{18}$, $R_{20}$, $R_{22}$ and $(R_{30})^{ijk}$ are genuine relations; $R_{24}^{(0,3)}$,
$(R_{26}^{(0,4)})^i$ and the singlet $R_{30}^{(0,5)}$ are not, they are graviton classes, and
$(R_{26}^{(0,4)})^i=-\tfrac13Q^i_+R_{24}^{(0,3)}$ modulo exact terms (\R{09}).
\item Table~1: reproduced with the corrected factor through $j=40$ (signed). The $[0,0]$ entries at $j=24$, 30 are
gravitons; the $-t^{30}[3,0]$ is the genuine class of \R{07}; the other entries through $j=40$ are $P$-odd; $[0,0]$ at
$j=36$ and $[1,0]$ at $j=38$ are undetermined (section~\ref{sec:su3j24}).
\end{itemize}
\item[Lee, thesis \cite{LeeThesis}.] Restates the diagonal criterion (same remark as above); eq.~6.33 has the correct
descendant factor.
\item[de Mello Koch, Kim, Kim, Lee, Lee \cite{dMK}.] Eqs.~(2.24)--(2.25): with three corrections to the printed formulas
(section~\ref{sec:su3j30}) the operator $O^3_{30}$ is closed at $N=3$ and its $x_1^3$ coefficient is
$-\tfrac{648}{5}O^2_{24}$, as stated. It is the first $\SU(3)$ black hole (\R{07}, \R{16}); their next member
$O^3_{32}$ is the $(5;8,4,4)$ class found by our $j=32$ scan. The paper regards $j=24$ as the $\SU(3)$ threshold
(following \cite{CCKLL}); there is no class there. Their picture of fused classes is confirmed below $j=34$, where
every $\SU(3)$ black hole is fused, and the first unfused core appears at $j=34$ (\R{24}).
\item[Gadde, Lee, Raj, Tomar \cite{GLRT}.] Agrees: the closed forms of the $\SO(7)$ and $\USp(6)$ BMN indices (through
$j=18$, as far as their printed numerators determine them) and of the difference through $j=30$; the singlet at
$j=18$; the numbers of sec.~4.3 (903 / 559 / 124 against 123); the representative eq.~(4.14) is closed on
$\mathfrak{so}(7)$ and not on $\mathfrak{so}(9)$. New relative to it: the $\SO(9)$, $\SO(11)$ classes and the index
differences for $N=4,5,6$ (\R{04}, \R{06}, \R{11}, \R{21}); the restriction to $\Gtwo$ (\R{10}).
\item[Chang, Lin \cite{CLSduality}.] Agrees: Table~1 in full, both columns and both groups, with the
monotone/fortuitous split (BMN column 903 / 559 / 124 against 123; chiral-ring column 826 / 741 / 85 against
825 / 741 / 84), and the statement that the extra $\SO(7)$ chiral-ring graviton is invisible on the Coulomb branch
(\R{05}, \R{25}).
\item[Choi, Lee \cite{ChoiLee}.] Agrees: the one-loop lifting of the $\SO(7)$ class, with the exact-space rank 741 in
the target sector, and $Q_1\mathrm{tr}(\ph_1\ps_1)=4(n-2)\mathrm{tr}(\lambda_+\lambda_-)$ for $\mathfrak{so}(n)$ (\R{20}).
Extended: the $\SO(9)$ and $\SO(11)$ classes are lifted too, through their $\SO(7)$ core.
\item[Kim, Kim, Lee, Park \cite{KKLP}.] Agrees: every one-loop verdict in our common range ($O_0$, $w_2O_0$, $O_1'$
unlifted; $O_1$ lifted, exact-space rank 912 at $j=36$); their eq.~(2.19) with an explicit $Q_2$; the entropy excess
1.21\% and the $\SO(7)$ ratio 1.0121. Extended: the complete $\SU(2)$ one-loop table through $j=44$; for $O_0$, whose higher
orders the paper leaves open (p.~6), no admissible partner at any order given the selection rule, $d_2[O_0]=0$
directly, and protection at every coupling by the index (section~\ref{sec:su2}); the charge
dependence of the entropy excess (section~\ref{sec:entropy}).
\item[Budzik, Kulp \cite{BK}.] We use their one-loop supercharge $Q_1$ (with an $\mathfrak{so}(n)$ completeness term for
the orthogonal groups) and their nilpotency relation eq.~(12). The two-loop supercharge on derivative-free words with
$K_2=-\tfrac16$ (section~\ref{sec:q2}) is, as far as we know, new; \cite{BGKWWY} state for $\cN=1$ theories that
$Q_1^2\ne0$ forces $Q_2\ne0$ and that a single two-loop diagram contributes, consistent with our counting.
\item[Kinney, Maldacena, Minwalla, Raju \cite{KMMR}.] Sec.~6.1: the eigenvalue count of finite-$N$ gravitons misses
the non-Coulomb gravitons (section~\ref{sec:noncoulomb}). For the gaugino-free part of the $\frac18$-BPS chiral ring
that our BMN data cover, their statements are not contradicted in our range.
\end{description}

\section{Cross-checks}\label{sec:checks}

\begin{longtable}{@{}>{\raggedright\arraybackslash}p{0.27\textwidth}>{\raggedright\arraybackslash}p{0.68\textwidth}@{}}
\toprule
check & where applied, and outcome\\
\midrule
\endhead
Exact-dimension self-check & Every sector of every highest-weight scan (all groups, BMN and full letters): rank of the
highest-weight vectors at the functionals $=$ exact multiplicity, per fermion degree.\\
Euler characteristic $=$ exact index & Every column of every complete scan ($\SU(2)$ 446/446; $\SU(3)$ 172/172 at
$j\le30$, 30/30 at $j=32$, the $(8,5,4)$ column; $\SU(4)$ 145/145 and 27/27; $\SO(5)$, $\SO(7)$, $\SO(9)$,
$\SO(11)$, $\USp(6)$, $\Gtwo$ all columns).\\
Completeness by enumeration & Expected (column, $2J$) lists compared with the data files for all complete scans.\\
Second random seed & Every positive (non-graviton or non-Coulomb) count, and the one-loop tests (two seeds each); the
$\SU(3)$ core column at $2J=4$--$7$; identical results throughout.\\
Larger generator sets & $\SU(3)$ $j=30$ with $n\le4,5$; $\SO(7)$ with $n\le10$; unchanged counts.\\
Independent engines & Weight-space direct cohomology (\file{bmn.symdirect}) against the highest-weight engine
(\file{hwcoh}) on $\SU(2)$ $j=16$, $\SU(3)$ $j=14$--$18$, $\SO(9)$ $j=22$; full-letter engine (\file{full}) against
\file{hwcoh} on all derivative-free sectors at levels 18, 20; the $\SU(2)$ non-graviton counts by the one-loop engine;
$\SO(5)$ with $\USp(4)$ letters.\\
Literature anchors & Table~2 of \cite{CCKLL}; eq.~(3.4) of \cite{CCKLL}; Table~1 of \cite{CLSduality}; sec.~4.3 and
the index closed forms of \cite{GLRT}; rank 741 of \cite{ChoiLee}; rank 912 and the one-loop verdicts of \cite{KKLP};
the 1.21\% of \cite{KKLP}; $-\tfrac{648}{5}$ of \cite{dMK}.\\
Exact arithmetic & Rational preimage of eq.~(4.8) from five primes (CRT); diagonal-function counts as exact linear
algebra; the $\SU(2)$ diagonal dimensions over $\mathbb Q$; exact multiplicities by Haar integration (rounding
errors below $2\times10^{-7}$, results integral).\\
Controls & Exact controls $Q_0u$ in every rank test; known-exact operators reported exact (e.g.\ $O_1^{(0,3)}$ of
\cite{CCKLL}); planted non-exact rows (``$z+$noise'') reported non-exact; negative controls for $K_2$.\\
Rank stability & Doubling, and for the one-loop lifts tripling, of the evaluation points; spanning certificates.\\
Symmetries & $(2J_1,2J_2)\leftrightarrow(2J_2,2J_1)$ in every full-letter sector; $\SU(3)_R$ covariance of the block
restriction coefficients ($1:-2:1$); Weyl symmetry of the diagonal dimensions.\\
S-duality & Equal Coulomb-type graviton counts for $\SO(2N+1)$ and $\USp(2N)$ in all compared sectors.\\
Index predictions & Kernel parity forecasts confirmed: $\SU(2)$ in all 446 columns; $\SU(3)$ at $j=32$ and in
$(8,5,4)$.\\
A-priori predictions & The $\SU(2)$ one-loop fate table predicted from selection rules before the direct computation
(\T{37}); 16 of 16 sectors agree.\\
Independent review & The results \R{01}--\R{18} were reviewed by a separate set of AI agents, with a cross-review of the
claim list by an AI model from another developer (2026-10-01/02): the prose was consistent with the stored data and no
stored number was wrong; the defects found (wording, claimed rigour, reproducibility) were corrected, including the
status of the kernel parity index (section~\ref{sec:parity}) and the energy bound $E\ge j/3$. Reproducer scripts:
\file{scripts/review/independent-checks}.\\
Library tests & \file{code/tests}, run when this repository was assembled: 82 passed, 4 skipped (three slow
tests that run with \texttt{FORTUITY\_SLOW=1}, and one that needs \texttt{fork}, disabled on macOS).\\
\bottomrule
\end{longtable}

What random evaluation proves, in one sentence: computed ranks never exceed true ranks over $\mathbb F_p$, so
certified null results are exact over $\mathbb F_p$, and a bad prime could only add pairs of classes.

\section{Open questions}\label{sec:open}

\begin{enumerate}
\item $d_2$ of the $\SU(3)$ class R-007 (the only open order without the Coulomb-branch assumption); the validity of
the superconformal selection rule beyond one loop; the lowest-energy status of R-007 (sectors $j=34$, $2J\ge10$ and
$j=36$, $2J=12$).
\item The $\SU(2)$ survivors at higher order: order 3 of $w_2O_0$ (three non-Coulomb classes in the target) and of
$O_1'$; $d_2$ beyond the multiplet of $O_0$.
\item $\SU(3)$ cores: the column $(8,6,4)$ at $j=36$ (fused capacity 0); the one-loop fate of the $j=34$ core; the
$P$-even entries $[0,0]$ at $j=36$ and $[1,0]$ at $j=38$ of Table~1 of \cite{CCKLL} (a direct computation of the
singlet column at $j=36$ was started and not finished).
\item The $\SU(4)$ threshold ($j\ge32$); $\SU(5)$ by direct cohomology.
\item $\SO(5)$ at $j=30$; direct $\SO(13)$ cohomology at $j=30$ (three fermionic classes predicted); $\SO(7)$ and
$\USp(6)$ at level 30 (sizing done, scan not); $\USp(6)$ at $j=26$; the 13 unclassified full-letter multiplets at
level 22.
\item The $C$-twisted index with all letters at the charges of R-007, which is coupling-independent and could protect a
$C$-odd state non-perturbatively.
\item Are non-Coulomb gravitons ever lifted as targets in $\SU(N)$? Which fortuitous classes above the thresholds
survive?
\item $\SO(2N)$ groups (Pfaffian invariants), and the $\USp(2N)$ thresholds beyond $j=28$.
\end{enumerate}

\paragraph{Note added.} While this summary was being prepared we learned of the independent work of Kim, Lee, Lee and
Song \cite{KLLS}, whose results overlap with ours in $\SU(3)$ and $\SU(4)$. We thank them for sharing their results
and for coordinating the posting.

\appendix

\section{Index of results, code and data}\label{app:index}

Paths are relative to the root of this repository. \file{code/fortuity} is the library (with tests in
\file{code/tests}); \file{scripts/<area>/<project>} hold the drivers, cluster job files and run outputs behind each
result; \file{data} holds the results in plain formats, each file with a \file{.meta.json} stating what it contains,
the code and command that produced it and how it was checked. ``Labels'' are those of the collaboration's records,
which the code comments use.

\small
\begin{longtable}{@{}>{\raggedright\arraybackslash}p{0.07\textwidth}>{\raggedright\arraybackslash}p{0.25\textwidth}>{\raggedright\arraybackslash}p{0.62\textwidth}@{}}
\toprule
label & content & code, data, checks\\
\midrule
\endhead
\R{01} & $\SU(2)$ threshold $j=24$, descendants, descendant factor & \file{code/fortuity/bmn/direct.py},
\file{fullindex.py}, \file{gravindex.py}, \file{run_su2_direct.py}; \file{scripts/bmn/index-and-direct-cohomology}
(\file{run_index.py}, \file{extract_core.py}, \file{check_su2_index.py}, \file{test_threshold_ops.py});
\file{data/su2/bmn_direct_cohomology_j28.json}, \file{bmn_index_j40.json}; check
\file{python3 -m fortuity.bmn.run_su2_direct 28 OUT.json --check data/su2/bmn_direct_cohomology_j28.json};
\file{code/tests/test_result_drivers.py}.\\
\R{02} & $\SU(3)$ $j=24$: 29 multi-graviton classes; criterion false & \file{code/fortuity/bmn/symdirect.py},
\file{r002_check.py}; \file{scripts/bmn/index-and-direct-cohomology} (\file{run_sym.py},
\file{test_diag_criterion.py 3,4,4,4}, \file{test_grav_sym.py}, \file{test_grav_robust.py},
\file{explicit_graviton_class.py}, \file{verify_explicit_class.py});
\file{data/su3/bmn_table2_weyl_symmetric_j24.json}, \file{graviton_class_vanishing_on_diagonal_j24.json}.\\
\R{03} & $\SU(4)$, $\SU(5)$ diagonal-count index & \file{code/fortuity/bmn/fullindex.py}, \file{gravindex.py};
\file{scripts/bmn/index-and-direct-cohomology} (\file{run_index.py}, \file{run_su4_energy.py});
\file{data/su4/bmn_index_j30.json}, \file{j28_highest_weight_sectors.json}, \file{data/su5/bmn_index_j28.json}.\\
\R{04} & $\SO(2N+1)$ vs $\USp(2N)$ indices, $N=3,4,5$ & \file{code/fortuity/bmn/groupindex.py};
\file{data/so_sp/so7_sp3_bmn_index.json}, \file{so9_sp4_bmn_index.json}, \file{so11_sp5_bmn_index.json};
\file{code/tests/test_groupindex.py}.\\
\R{05} & $\SO(7)$/$\USp(6)$ at $(1;3,3,3)$ & \file{code/fortuity/bmn/symdirect.py}, \file{evaluate.py};
\file{data/so_sp/so7_direct_cohomology_1333.json}, \file{usp6_direct_cohomology_1333.json}.\\
\R{06} & $\SO(9)$ class at $(1;5,3,3)$ & as \R{05}; \file{data/so_sp/so9_direct_cohomology_533.json},
\file{usp8_direct_cohomology_533.json}, \file{so9_hw_cohomology_j22.jsonl},
\file{so9_hw_cohomology_j22_seed2_533.jsonl}; closure test
\file{scripts/bmn/so-usp-g2-scans/extensions/closure_check.py} (output \file{closure_check.txt}).\\
\R{07} & $\SU(3)$ complete scan $j\le30$; class at $j=30$ & \file{code/fortuity/hwcoh}; cluster jobs
\file{scripts/bmn/hw-scans-cluster}; \file{data/su3/bmn_hw_cohomology_j30.json}; representative (\T{21})
\file{scripts/bmn/su3-level30-representative/representative.py} and \file{runs/}.\\
\R{08} & $\SU(4)$ $j\le28$ & \file{code/fortuity/hwcoh}; \file{scripts/bmn/hw-scans-cluster};
\file{data/su4/bmn_hw_cohomology_j28.json}.\\
\R{09} & Eq.~(4.8) of \cite{CCKLL} exact; (A.4); diagonal mechanism; Table~1 & \file{scripts/bmn/su3-level24}
(\file{build_O.py}, \file{build_O6.py}, \file{singlet_preimage*.py}, \file{hw_relation_test.py},
\file{triple_check.py}, \file{decompose*.py}, \file{full_complex_check.py}, \file{f_restricted_exactness.py},
\file{R24_check.py}, \file{descendant_check.py}, \file{diag_vanishing.py}, \file{table1_core.py},
\file{singlet_column_j36.py}); \file{data/su3/cckll_eq4p8_operator_j24.json}, \file{cckll_eq4p8_preimage_j24.json},
\file{cckll_A4_relations_exactness.json}; \file{code/tests/test_r009_preimage.py}.\\
\R{10} & $\Gtwo$ completeness to $j=18$; restriction of the $\SO(7)$ class & \file{code/fortuity/bmn/g2_checks.py}
(\file{python3 -m fortuity.bmn.g2_checks restriction}, \file{completeness 18});
\file{data/g2/g2_so7_restriction_1333.json}; \file{code/tests/test_g2_checks.py}.\\
\R{11} & $\SO(11)$/$\USp(10)$ at $j=26$ & \file{code/fortuity/hwcoh} (groups \file{SO_odd}, \file{Sp}); job files
\file{scripts/bmn/so-usp-g2-scans/cluster/t020_*.sbatch}; \file{data/so_sp/so11_hw_cohomology_j26*.jsonl},
\file{usp10_hw_cohomology_j26*.jsonl}.\\
\R{12} & $\SO(5)$ none for $j\le28$ & \file{code/fortuity/hwcoh}; \file{scripts/bmn/so5-scan} (\file{run_so5.sh},
\file{check_so5_driver.py}); \file{scripts/bmn/hw-scans-cluster} (\file{hw_par.sbatch}, \file{so5_json.py});
\file{data/so5/bmn_hw_cohomology_j28.json} (and the $j\le24$, $j\le26$ snapshots).\\
\R{13} & $\SO(7)$, $\Gtwo$, $\USp(6)$ at $j=20$--$26$ & \file{code/fortuity/hwcoh} (\file{restrict.py},
\file{evalg.py}); \file{scripts/bmn/so-usp-g2-scans} (\file{restrict_scan.py}, \file{euler_check.py},
\file{extra_checks.py}, \file{cluster/t023_*.sbatch}); \file{data/g2} (files \file{*_j20*}, \file{*_j22_j24*},
\file{*_j26*}, \file{so7_robustness_seed2_gen10.jsonl}, \file{extra_564448}).\\
\R{14} & $\SU(4)$ $j=30$ & \file{code/fortuity/hwcoh}; \file{scripts/bmn/hw-scans-cluster/cluster/launch_su4_j30.sbatch};
\file{data/su4/bmn_hw_cohomology_j30.json}.\\
\R{15} & Non-Coulomb gravitons & \file{code/fortuity/hwcoh/noncoulomb.py}
(\file{python -m fortuity.hwcoh.noncoulomb GROUP N GEN_NMAX OUT DATA...}); \file{data/noncoulomb}
(\file{README.meta.json}); \file{code/tests/test_noncoulomb.py}.\\
\R{16} & Block restriction, fused classes, the operator of \cite{dMK} & \file{code/fortuity/bmn/blockrestrict.py},
\file{classify.py}, \file{dmk.py}, \file{code/fortuity/hwcoh/classreps.py}; \file{scripts/bmn/block-restriction/code}
(\file{run_su3_to_su2.py}, \file{run_so_tower.py}, \file{make_r016_data.py}) and \file{runs/};
\file{data/su3/r007_block_restriction_su2xu1.json}, \file{r007_x1cubed_image_su2.json}, \file{dmk_o330_hw.json},
\file{data/so_sp/*block_restriction*.json}, \file{*nongraviton_rep*.json}; \file{code/tests/test_blockrestrict.py},
\file{test_dmk.py}.\\
\R{17} & Parity lemma, twisted and kernel parity indices & \file{code/fortuity/bmn/twistedindex.py}, \file{cparity.py},
\file{gravindex.py}; \file{scripts/index/c-parity/code} (\file{twisted_index.py}, \file{noncoulomb_parity.py},
\file{su2_predictions.py}, \file{diag_counts.py}, \file{check_r007_cparity.py}, \file{atlas.py},
\file{su4_twisted_scan.py}) and \file{runs/}; the review of its rigour in \file{scripts/index/c-parity/c19};
\file{data/su3/twisted_index_j30.json}, \file{data/su4/twisted_index_j30.json}, \file{data/su3/kernel_parity_j40.json},
\file{data/su2/kernel_parity_j44.json}; \file{code/tests/test_twistedindex.py}.\\
\R{18} & $\SO(7)$, $\SO(9)$, $\SO(11)$ lowest levels certified & \file{code/fortuity/hwcoh};
\file{scripts/bmn/so-odd-thresholds} (\file{make_tasks.py}, \file{dims.sbatch}, \file{col.sbatch}, \file{euler.sbatch},
audit \file{analyze.py}); \file{data/so_sp/so7_hw_cohomology_j18.jsonl}, \file{so9_hw_cohomology_j20.jsonl},
\file{so9_hw_cohomology_j22.jsonl}, \file{so11_hw_cohomology_j24.jsonl},
\file{so9_hw_cohomology_533_j22_control.jsonl}.\\
\R{19} & $\SO(7)$, $\USp(6)$, $\Gtwo$ at $j=28$ & \file{code/fortuity/hwcoh}; \file{scripts/bmn/so-usp-g2-scans}
(\file{make_sector_tasks.py}, \file{extra_checks.py}, \file{cluster/t023_scan_v2.sbatch},
\file{cluster/t023_j28_checks.sbatch}); \file{data/g2/*_j28.jsonl}, \file{data/g2/checks_j28}.\\
\R{20} & One-loop lifting of the $\SO(2N+1)$ classes & \file{code/fortuity/full/q1.py}, \file{lifting.py}
(\file{python -m fortuity.full.lifting so7|so9|so11}); \file{scripts/oneloop/so-odd-classes} (development scripts and
\file{runs/}); \file{code/tests/test_q1_groups.py}, \file{test_lifting.py}.\\
\R{21} & $\SO(13)$ vs $\USp(12)$ index & \file{code/fortuity/bmn/groupindex.py}, \file{fugacityindex.py};
\file{data/so_sp/so13_sp6_bmn_index.json}; log \file{scripts/bmn/so-usp-g2-scans/extensions/n6_j30_fugacity.txt};
\file{code/tests/test_fugacityindex.py}.\\
\R{22} & $\Gtwo$ $j\le18$ & \file{code/fortuity/hwcoh} (group \file{G2}); \file{data/g2/g2_hw_cohomology_j18.jsonl};
column list and log \file{scripts/bmn/so-usp-g2-scans/extensions/g2}.\\
\R{23} & Non-Coulomb gravitons, $\SO$ series and $j=28$ & \file{code/fortuity/hwcoh/noncoulomb.py};
\file{data/noncoulomb/noncoulomb_*_j18|j22|j24|j28.jsonl}, \file{so_odd_j28_extension.meta.json},
\file{logs_2026-10-06}.\\
\R{24} & $\SU(3)$ core at $j=34$ & \file{code/fortuity/hwcoh} (\file{run_sectors.py}, \file{cores.py},
\file{classreps_par.py}); \file{scripts/bmn/su3-cores} (\file{cluster/sector.sbatch}, \file{analysis/make_core_data.py},
\file{reps/run_reps.py}, \file{reps/reps.sbatch}, \file{reps/make_reps_data.py}) and its \file{results*};
\file{data/su3/hw_cohomology_8_5_4.json}, \file{core_restrictions_j30_j34.json}; \file{code/tests/test_cores.py},
\file{test_classreps_par.py}.\\
\R{25} & Full letters, $\SO(7)$/$\USp(6)$ levels 18--22 & \file{code/fortuity/full} (\file{sector.py}, \file{gdims.py},
\file{levels.py}, \file{gravitons.py}); \file{scripts/full-letter/so7-usp6} (\file{enum_level.py},
\file{sector_group.py}, \file{hwgrav.py}, \file{runs/*.py}); \file{data/fullso} (\file{fullso.meta.json});
\file{code/tests/test_full_groups.py}.\\
\midrule
\T{19} & One loop for $\SO$ classes (= \R{20}) & see \R{20}.\\
\T{21} & Explicit representative of R-007 & see \R{07}.\\
\T{25} & $\SO$ thresholds (= \R{18}) & see \R{18}.\\
\T{27} & $\SU(2)$ complete scan $j\le44$ & \file{code/fortuity/hwcoh}; \file{scripts/bmn/su2-level44-scan}
(\file{setup.sbatch}, \file{su2_scan.sbatch}, \file{su2_col.sbatch}, \file{seed_columns.py}, \file{watch_su2*.sbatch},
\file{tasks_su2_j44.txt}); per-sector records \file{data/su2/hw_cohomology_j44_records}, assembled
\file{data/su2/bmn_hw_cohomology_j44.json} (\file{python -m fortuity.hwcoh.to_json}); non-Coulomb test
\file{data/noncoulomb/noncoulomb_su2_j44.jsonl}.\\
\T{30}, \T{31}, \T{40} & One loop, all orders and two loops for the $\SU(2)$ and $\SU(3)$ classes &
\file{scripts/oneloop/su-threshold-classes}: \file{oneloop_test.py}, \file{fullpar.py}, \file{gpu_rank.py},
\file{cpu_rank.py} (one loop); \file{c_parity.py}, \file{twisted_index.py} ($C$); \file{size_r007_orders.py},
\file{membership.py}, \file{hwcheck.py}, \file{abelian_lw*.py}, \file{su2_lw.py}, \file{test_q1_symm.py} (orders 2--4,
selection rule); \file{su2_index_tower.py} (index); \file{q2bmn.py}, \file{test_q2.py}, \file{test_q2_tracerel.py},
\file{d2_make.py}, \file{d2_report.py} (two loops); \file{cluster/} (job files); \file{runs/} (outputs, including
\file{t040_oneloop_table.json}, \file{t040_twoloop_table.json}, \file{r007_*}).\\
\T{32} & $\SU(3)$ cores (= \R{24}) & see \R{24}.\\
\T{33} & Restriction and all-orders protection & \file{scripts/oneloop/protection/code} (\file{so9_intertwining.py},
\file{so11_intertwining.py}, \file{su2_channels.py}) and \file{runs/}.\\
\T{34} & Kernel parity for $\SO$, $\USp$, $\Gtwo$ & \file{scripts/index/kernel-parity-so-usp-g2/kernel_index.py} and its
outputs.\\
\T{37} & A-priori $\SU(2)$ fate table & \file{scripts/oneloop/su2-fate-predictions/scripts} (\file{run_all.sh}) and
\file{outputs/}.\\
\T{38} & Bonus index with parity & \file{scripts/entropy/bonus-index-parity} (\file{cardy_lib.py},
\file{cardy_validate.py}, \file{cardy_kklp.py}, \dots, outputs \file{*.out}).\\
\T{41} & Charge dependence of the lifting excess & \file{scripts/entropy/lifting-excess} (\file{code/},
\file{physics/}, independent implementation \file{indep/}, \file{cluster/}, \file{outputs/}, \file{figures/}).\\
\T{42} & $\SU(3)$ complete scan $j=32$ & \file{code/fortuity/hwcoh}; \file{scripts/bmn/su3-level32-scan}
(\file{dims.sbatch}, \file{col.sbatch}, task lists); \file{data/su3/hw_cohomology_j32_records}; Euler check
\file{python -m fortuity.hwcoh.to_json 'data/su3/hw_cohomology_j32_records/*.jsonl' 3 OUT.json}.\\
\T{43} & $\SO(13)$ index, $\Gtwo$ below $j=20$, $\SO(9)$ closure & \file{scripts/bmn/so-usp-g2-scans/extensions}.\\
\T{44}--\T{48} & $\SO(9)$ level 22, corrections, non-Coulomb extension, library promotion & see \R{18}, \R{19},
\R{23}, \R{25}.\\
other & labels in code comments & \T{20} (computations of \R{04}, \R{11}), \T{23} (\R{10}, \R{13}, \R{19}, \R{22}),
\T{24} (\R{12}), \T{28} (\R{16}), \T{29} (\R{17}); \T{16}, \T{39}: promotion of the full-letter code and of the
$\mathfrak{so}(n)$ one-loop supercharge into \file{code/fortuity/full}; \T{01}, \T{10}: conventions and the
operational definition (section~\ref{sec:method}).\\
\midrule
\textsf{D-005} & Operational graviton definition & section~\ref{sec:method}.\\
\textsf{D-006} & BMN conventions & section~\ref{sec:setup}; \file{code/fortuity/bmn/words.py}.\\
other & Validation of the full-letter code & \file{scripts/full-letter/su-checks} (\file{validate_su2.py},
\file{validate_hw_su2.py}, \file{run_bmn_check.py}, \file{classes_full.py}, \file{test_q1.py}).\\
other & Independent review & \file{scripts/review/independent-checks} (\file{reproducers/}, \file{followup/}).\\
\bottomrule
\end{longtable}
\normalsize

\section{Reproduction}\label{app:repro}

\paragraph{Installation.} Python $\ge3.10$ with numpy, scipy, sympy (\file{code/pyproject.toml}):
\begin{small}
\begin{verbatim}
python3 -m pip install -e 'code[dev]'   # or: export PYTHONPATH=$PWD/code
cd code && python3 -m pytest -q         # FORTUITY_SLOW=1 also runs the slow tests
\end{verbatim}
\end{small}
The highest-weight engine runs on numpy alone; the optional FLINT/BLAS kernels are built with the scripts in
\file{code/fortuity/hwcoh/flint} (FLINT 3, OpenBLAS or MKL). The GPU one-loop ranks need CuPy.

\paragraph{Quick checks} (from the repository root unless stated; times on an Apple-silicon laptop).
\begin{small}
\begin{verbatim}
cd code
python3 -m fortuity.bmn.r002_check              # R-002 operator, seconds
python3 -m fortuity.bmn.g2_checks restriction   # R-010, ~20 s
python3 -m fortuity.full.lifting so7            # R-020, SO(7) lifting
python3 -m fortuity.bmn.run_su2_direct 28 /tmp/su2.json \
        --check ../data/su2/bmn_direct_cohomology_j28.json   # R-001, ~8 min
cd ..
PYTHONPATH=code python3 -m fortuity.hwcoh.to_json \
        'data/su3/hw_cohomology_j32_records/*.jsonl' 3 out.json   # T-042 Euler check
export PYTHONPATH=$PWD/code REPO=$PWD
python3 scripts/bmn/su3-level24/table1_core.py   # Table 1 of CCKLL, seconds
cd scripts/bmn/su3-level24
python3 diag_vanishing.py       # diagonal mechanism, seconds
python3 descendant_check.py 1   # R_26 = -1/3 Q_+ R_24 + exact, ~1 min
cd ../so-odd-thresholds
python3 analyze.py 5 24 ../../../data/so_sp/so11_hw_cohomology_j24.jsonl
cd ../../oneloop/su-threshold-classes
python3 test_q2.py 12           # K_2 = -1/6, seconds
python3 su2_index_tower.py      # index tower of O_0, ~10 min
cd ../../entropy/bonus-index-parity && python3 cardy_validate.py
cd ../../full-letter/su-checks && python3 validate_su2.py 6
\end{verbatim}
\end{small}
Scripts that are not modules need \file{PYTHONPATH} to contain \file{code}; scripts under
\file{scripts/oneloop/su2-fate-predictions} also read the repository root from \file{REPO}.

\paragraph{Cluster job files.} The \file{*.sbatch} and \file{env.sh} files are kept as they were run, with
machine-specific values replaced by placeholders: \var{REPO} (the copy of this code used on the cluster),
\var{SCRATCH}, \var{STORE}, \var{VENV} (scratch, storage and virtual-environment locations), \var{ACCOUNT},
\var{USER}, \var{EMAIL}, \var{HOME}, \var{FLINT\_PREFIX}. In the layout used there the library sat under
\var{REPO}\texttt{/code} and each job wrote per-sector records into a \file{results} directory; those records are
the files in \file{data}. The largest scans took of order $10^2$--$10^3$ core-hours ($\SU(3)$ to $j=30$
about 400, $\SU(4)$ to $j=28$ about 900 and at $j=30$ about 1500); the one-loop and higher-order $\SU(3)$ tests needed
nodes with up to 0.5\,TB of memory or an H200 GPU.

\begin{thebibliography}{99}\small

\bibitem{CLWords} C.-M. Chang and Y.-H. Lin, \emph{Words to describe a black hole}, JHEP \textbf{02} (2023) 109,
\arxiv{2209.06728}.

\bibitem{CLCovering} C.-M. Chang and Y.-H. Lin, \emph{Holographic covering and the fortuity of black holes},
\arxiv{2402.10129}.

\bibitem{CKLP} S. Choi, S. Kim, E. Lee and J. Park, \emph{The shape of non-graviton operators for $SU(2)$},
\arxiv{2209.12696}.

\bibitem{CCKLL} J. Choi, S. Choi, S. Kim, J. Lee and S. Lee, \emph{Finite $N$ black hole cohomologies},
\arxiv{2312.16443}.

\bibitem{LeeThesis} S. Lee, \emph{Quantum anatomy of supersymmetric black holes in AdS spacetimes}, PhD thesis
(2024), \arxiv{2409.03862}.

\bibitem{dMK} R. de Mello Koch, M. Kim, S. Kim, J. Lee and S. Lee, \emph{Brane-fused black hole operators},
\arxiv{2412.08695}.

\bibitem{GLRT} A. Gadde, E. Lee, R. Raj and S. Tomar, \emph{Probing non-graviton spectra in $\cN=4$ SYM via BMN
truncation and S-duality}, \arxiv{2506.13887}.

\bibitem{CLSduality} C.-M. Chang and Y.-H. Lin, \emph{Violation of S-duality in classical $Q$-cohomology},
\arxiv{2510.24008}.

\bibitem{ChoiLee} J. Choi and E. Lee, \emph{Konishi lifts a black hole}, \arxiv{2511.09519}.

\bibitem{BK} K. Budzik and J. Kulp, \emph{Loop corrected supercharges from holomorphic anomalies},
\arxiv{2512.07771}.

\bibitem{KKLP} S. Kim, S. Kim, S. Lee and J. Park, \emph{Quantum black hole cohomologies}, \arxiv{2606.27955}.

\bibitem{KLLS} S. Kim, J. Lee, S. Lee and J. Song, \emph{Black hole cohomologies in $\cN=4$ SYM},
\arxiv{2610.20340}.

\bibitem{BGKWWY} K. Budzik, D. Gaiotto, J. Kulp, B. R. Williams, J. Wu and M. Yu, \emph{Semi-chiral operators in 4d
$\cN=1$ gauge theories}, JHEP \textbf{05} (2024) 245, \arxiv{2306.01039}.

\bibitem{KMMR} J. Kinney, J. Maldacena, S. Minwalla and S. Raju, \emph{An index for 4 dimensional super conformal
theories}, Commun.\ Math.\ Phys.\ \textbf{275} (2007) 209, \arxiv{hep-th/0510251}.

\bibitem{Schwartz} J. T. Schwartz, \emph{Fast probabilistic algorithms for verification of polynomial identities},
J.\ ACM \textbf{27} (1980) 701.

\bibitem{Zippel} R. Zippel, \emph{Probabilistic algorithms for sparse polynomials}, in \emph{EUROSAM '79}, Lecture Notes
in Computer Science \textbf{72} (1979) 216.

\bibitem{vMS} A. von Manteuffel and R. M. Schabinger, \emph{A novel approach to integration by parts reduction},
Phys.\ Lett.\ B \textbf{744} (2015) 101, \arxiv{1406.4513}.

\bibitem{Peraro} T. Peraro, \emph{Scattering amplitudes over finite fields and multivariate functional
reconstruction}, JHEP \textbf{12} (2016) 030, \arxiv{1608.01902}.

\end{thebibliography}

\end{document}
